\chapter{Lagerberg superforms}
\index{superform}\index{Lagerberg differential}
\label{chap:superforms}

The local differential calculus of Chambert--Loir and Ducros begins on real
affine space. Its two differentials are not two derivatives in different
directions: they are two labelled copies of the same real derivative. This
simple device retains the bidegrees, positivity, and Hessian calculus familiar
from complex geometry.

Throughout this chapter $V$ is an $n$-dimensional real vector space and
$U\subset V$ is open. We first work in affine coordinates and then prove that
the definitions do not depend on them.

It is worth being precise at once about what the doubling device does and
does not accomplish, because the analogy with complex geometry is close
enough to be dangerous.  On a complex manifold the decomposition
$\rd=\partial+\bar\partial$ is forced by the complex structure: the two
operators act on genuinely different objects, and a form of type $(1,0)$ is
not a form of type $(0,1)$.  On a real vector space there is no such
structure, and none is manufactured here.  What is manufactured is a
\emph{label}.  We take two copies of the cotangent space, agree to write the
first with a prime and the second with a double prime, and let one and the
same partial derivative $\partial/\partial x_i$ produce either
$\dprime x_i$ or $\ddprime x_i$ according to which copy we have decided to
feed.  Everything that follows---the bidegrees, the two differentials, the
involution exchanging them, the three positivity cones---is a consequence of
that bookkeeping and of nothing else.

The reward for the bookkeeping is that the Hessian of a smooth function
becomes a $(1,1)$-form,
\[
 \dprimeddprime u
 =\sum_{i,j}\frac{\partial^2u}{\partial x_i\partial x_j}\,
   \dprime x_i\wedge\ddprime x_j ,
\]
and its positivity becomes the positivity of a $(1,1)$-form in exactly the
sense used in complex pluripotential theory.  The whole of the local theory
in this part rests on this one identification: \emph{convex functions are the
plurisubharmonic functions of the real theory}.  The corresponding statement
is proved twice below, once for $C^2$ functions
(\Cref{thm:convex-ddc-positive}) and once in the distributional form
(\Cref{thm:ch09-distributional-convexity}) that the later chapters actually
use, because in the applications the potentials are piecewise affine and are
not twice differentiable anywhere along their break loci.

The price for the bookkeeping is that signs proliferate.  A superform of
bidegree $(p,q)$ has total degree $p+q$ and obeys the commutation rule of an
exterior algebra in that total degree; but the labels sit in two blocks, and
moving a double-primed factor past a primed one costs a sign that has nothing
to do with the bidegree of the form as a whole.  Almost every mistake in this
calculus is a mistake of that kind.  We therefore prove the sign rules rather
than assert them, and we test each one on an example in one or two variables
where the answer can be checked by hand.

\section{Orientation and prerequisites}
\label{sec:ch09-orientation}

Only three pieces of background are needed: multilinear algebra, the chain
rule, and the elementary Hessian criterion for convexity.  The two exterior
algebras below should be regarded as labelled copies.  The labels remember
which factor a differential occupies; they do not require a complex
structure on $V$.

By the end of the chapter the reader should be able to
\begin{enumerate}
\item compute $\dprime$, $\ddprime$, and $\dprimeddprime$ with all signs;
\item pull superforms back under affine maps and explain why nonlinear maps
  fall outside the functorial category;
\item recognize positivity in bidegree $(1,1)$ as positive semidefiniteness;
\item translate the Hessian test for convexity into positivity of
  $\dprimeddprime f$.
\end{enumerate}
The discussion is deliberately more explicit than the later analytic
chapters: every sign used in tropical integration and in the
Poincar\'e--Lelong formula is already visible on $\mathbf R^n$.

\subsection*{Roadmap}

\Cref{sec:ch09-algebra} constructs the bigraded algebra
$A^{\bullet,\bullet}(U)$, proves the wedge sign rules from the single
identity that the algebra is the exterior algebra of $V^*\oplus V^*$,
introduces $\dprime$, $\ddprime$ and the involution $J$, and proves the four
identities
\[
 (\dprime)^2=(\ddprime)^2=0,
 \qquad
 \dprimeddprime=-\ddprime\dprime,
 \qquad
 J\dprime J=\ddprime
\]
together with the Leibniz rules.  Two short calculations in one and in two
variables are carried out completely, and the resulting tables of basis
elements are the ones used in every later verification.

\Cref{sec:ch09-contraction} adds the two contraction operators
$\iota'_v$ and $\iota''_v$.  They are not needed for the algebra of forms,
but they are the operators in which the boundary traces of
\Cref{chap:tropical-integration} and the homogeneity conditions of
\Cref{sec:appD-supercurrents} are written, and their signs are of exactly the
kind that this chapter exists to pin down.

\Cref{sec:ch09-pullback} proves that the whole structure is functorial for
affine maps: pullback commutes with the wedge product, with $J$, and with
both differentials, whence the calculus is independent of the affine
coordinates used to define it.  The same computation shows why nonlinear maps
must be excluded, and this is the algebraic reason that the transition maps
of a tropical chart in \Cref{chap:tropical-spaces} are required to be
integral affine.

\Cref{sec:ch09-symmetry} sets up the dictionary between symmetric
$(p,p)$-forms and symmetric matrices indexed by $p$-element subsets, and
proves the complementary pairing formula.  This is the technical heart of the
chapter: once it is available, the three positivity notions of
\Cref{sec:ch09-positivity} become three familiar conditions on a symmetric
matrix, their inclusions and their duality become statements about convex
cones of matrices, and the standard counterexample separating them becomes a
computation with six coordinates that can be checked line by line.

\Cref{sec:ch09-convexity} identifies positivity of $\dprimeddprime u$ with
convexity of $u$ and computes the Monge--Amp\`ere form
$(\dprimeddprime u)^n$ in coordinates.
\Cref{sec:ch09-regularization} proves the regularization lemma---convolution
of a convex function with a radial kernel is convex, dominates the function,
and increases with the parameter---and deduces the distributional
characterization of convexity.  Both are used repeatedly from
\Cref{chap:cld-bedford-taylor} onwards.
\Cref{sec:ch09-dictionary} collects the dictionary and the sign checklist,
and the closing section separates Lagerberg's original formalism from the
Chambert--Loir--Ducros development.

\subsection*{What is deliberately absent}

Three things that a reader might expect are not in this chapter, and it is
better to say so now.

There is no integration theory.  A top-degree superform is not a measure
until an orientation and a Lebesgue normalization have been chosen, and the
choice is genuinely additional data: this is
\Cref{chap:tropical-integration}, where the calibration formalism is
introduced precisely because the naive Lagerberg integral is not invariant
under all changes of coordinates.  The present chapter fixes the orientation
$\tau_n$ and nothing more.

There is no sheaf theory and no cohomology.  Everything here happens on a
single open subset of a single real vector space.

There is no notion of a superform with singular coefficients.  The
coefficients are $C^\infty$ throughout, with the single exception of the
distributional statement in \Cref{sec:ch09-regularization}, where the object
being differentiated is a function and the derivatives are taken in the sense
of distributions on an open subset of $\mathbf R^n$.  Currents proper---the
objects that make $\dprimeddprime\max\{0,x\}$ legitimate---are constructed in
\Cref{chap:cld-currents}.

\section{The bigraded algebra}
\label{sec:ch09-algebra}

\begin{definition}
\label{def:ch09-superform-space}
The space of superforms of bidegree $(p,q)$ on $U$ is
\[
 A^{p,q}(U)=C^\infty(U)\otimes_{\mathbf R}
             \bigwedge^pV^*\otimes_{\mathbf R}\bigwedge^qV^*.
\]
In affine coordinates $x_1,\ldots,x_n$ we denote the two copies of $\rd x_i$ by
$\dprime x_i$ and $\ddprime x_i$. Thus
\[
 \alpha=\sum_{|I|=p,\ |J|=q}\alpha_{IJ}\,
 \dprime x_I\wedge \ddprime x_J.
\]
\end{definition}

Equivalently, a superform assigns to $x\in U$ a multilinear function
\[
 \alpha_x:(V^p)\times(V^q)\longrightarrow\mathbf R
\]
that is alternating in each block and depends smoothly on $x$.  This
description makes clear that a primed tangent vector and a double-primed
tangent vector both belong to $V$, but are inserted in different alternating
blocks.  In particular,
\[
 A^{p,q}(U)=0\quad\text{if}\quad p>n\ \text{or}\ q>n,
 \qquad
 A^{0,0}(U)=C^\infty(U).
\]

Three remarks on the definition are in order before any computation.

First, the two tensor factors $\bigwedge^pV^*$ and $\bigwedge^qV^*$ are the
same vector space.  Nothing distinguishes them except their position, and the
prime notation is exactly the device that records the position.  It is
legitimate, and often clarifying, to read $A^{p,q}(U)$ as the space of smooth
maps from $U$ to $\bigwedge^pV^*\otimes\bigwedge^qV^*$; the wedge product then
becomes an operation on the target and the differentials become operations on
the source.

Second, the definition is affine rather than linear in nature.  If $V$ is an
$n$-dimensional real \emph{affine} space with direction space
$\overrightarrow V$, one sets
\[
 A^{p,q}(U)
 =C^\infty(U)\otimes_{\mathbf R}
   \bigwedge\nolimits^p\overrightarrow V{}^*
   \otimes_{\mathbf R}
   \bigwedge\nolimits^q\overrightarrow V{}^*,
\]
and every statement below holds verbatim.  We suppress the arrow and treat a
choice of origin as part of the choice of coordinates; \Cref{sec:ch09-pullback}
proves that nothing depends on it.  The affine formulation is the one needed
in \Cref{chap:tropical-spaces}, where the charts of a tropical space are
open subsets of affine spaces and the transition maps are affine.

Third, the total dimension is
\[
 \dim_{\mathbf R}A^{p,q}(V)=\binom np\binom nq,
 \qquad
 \sum_{p,q}\dim_{\mathbf R}A^{p,q}(V)=4^n,
\]
which is $2^{2n}$, the dimension of the exterior algebra of a $2n$-dimensional
vector space.  That is not a coincidence, and it is the cleanest way to
organize the sign rules.

\begin{lemma}
\label{lem:ch09-exterior-of-double}
Let $W=V^*\oplus V^*$ and let $\iota',\iota'':V^*\to W$ be the two inclusions.
Then the assignment
\[
 \begin{aligned}
  \eta\otimes\xi\longmapsto{}&
  \iota'(\eta_1)\wedge\cdots\wedge\iota'(\eta_p)\\
  &{}\wedge\iota''(\xi_1)\wedge\cdots\wedge\iota''(\xi_q),\\
  &\eta=\eta_1\wedge\cdots\wedge\eta_p,
  \qquad \xi=\xi_1\wedge\cdots\wedge\xi_q,
 \end{aligned}
\]
extends to an isomorphism of graded algebras
\[
 \bigoplus_{p,q}A^{p,q}(V)
 \;\xrightarrow{\ \sim\ }\;
 \bigwedge\nolimits^{\bullet}W ,
\]
where the left side is graded by total degree $p+q$.  Under this
isomorphism $\dprime x_i$ and $\ddprime x_i$ correspond to
$\iota'(x_i)$ and $\iota''(x_i)$, and the elements
$\dprime x_I\wedge\ddprime x_J$, for $I$ and $J$ strictly increasing
multi-indices, form a basis.
\end{lemma}

\begin{proof}
For any two finite-dimensional vector spaces $W_1,W_2$ the exterior algebra
of the direct sum decomposes canonically as
\[
 \bigwedge\nolimits^{m}(W_1\oplus W_2)
 =\bigoplus_{p+q=m}
   \bigwedge\nolimits^{p}W_1\otimes\bigwedge\nolimits^{q}W_2 ,
\]
the isomorphism from right to left being induced by the multiplication map
$(\eta,\xi)\mapsto \eta\wedge\xi$ in $\bigwedge^\bullet(W_1\oplus W_2)$; that
this is an isomorphism is the standard computation with bases, and it is
clearly compatible with products up to the reordering sign recorded below.
Applying it with $W_1=W_2=V^*$ and summing over $m$ gives the displayed
bijection.  If $x_1,\ldots,x_n$ is a basis of $V^*$, then the $2n$ elements
$\iota'(x_1),\ldots,\iota'(x_n),\iota''(x_1),\ldots,\iota''(x_n)$ form a basis
of $W$, so the products of strictly increasing subfamilies form a basis of
$\bigwedge^\bullet W$; sorting each such product into its primed and
double-primed parts, which costs only a sign, produces exactly the family
$\dprime x_I\wedge\ddprime x_J$.
\end{proof}

The lemma is the whole content of the sign calculus: \emph{a superform is an
exterior form on a doubled space, and the bidegree merely records how many of
its arguments come from each copy}.  In particular the algebra is
graded-commutative for the total degree, and every sign question can be
answered by counting transpositions of degree-one factors.

For pure tensors the wedge product is fixed by
\[
(\eta\otimes\xi)\wedge(\eta'\otimes\xi')
 =(-1)^{q p'}(\eta\wedge\eta')\otimes(\xi\wedge\xi').
\]
Equivalently, all primed factors are placed before all double-primed factors.
It follows directly that
\[
             \alpha\wedge\beta=(-1)^{(p+q)(p'+q')}
             \beta\wedge\alpha.
\]

Both displayed formulas deserve a proof rather than an assertion, since they
are the source of every later sign.

\begin{proposition}
\label{prop:ch09-wedge-signs}
Let $\alpha\in A^{p,q}(U)$ and $\beta\in A^{p',q'}(U)$.
\begin{enumerate}
\item If $\alpha=\eta\otimes\xi$ and
  $\beta=\eta'\otimes\xi'$ are pure tensors, with
  $\eta\in\bigwedge^pV^*$, $\xi\in\bigwedge^qV^*$,
  $\eta'\in\bigwedge^{p'}V^*$, $\xi'\in\bigwedge^{q'}V^*$, then
  \[
   \alpha\wedge\beta=(-1)^{qp'}(\eta\wedge\eta')\otimes(\xi\wedge\xi').
  \]
\item In general
  $\alpha\wedge\beta=(-1)^{(p+q)(p'+q')}\beta\wedge\alpha$.
\item The product is associative and
  $C^\infty(U)$-bilinear, and $A^{\bullet,\bullet}(U)$ is a bigraded
  associative algebra with unit $1\in A^{0,0}(U)$.
\end{enumerate}
\end{proposition}

\begin{proof}
By \Cref{lem:ch09-exterior-of-double} we may compute inside
$\bigwedge^\bullet W$, where the product is the exterior product and
degree-one elements anticommute.  Write
$\alpha=\eta\wedge\xi$ and $\beta=\eta'\wedge\xi'$ with $\eta,\eta'$ products
of primed generators and $\xi,\xi'$ products of double-primed generators.
Then
\[
 \alpha\wedge\beta=\eta\wedge\xi\wedge\eta'\wedge\xi' .
\]
To restore the normal form, in which all primed factors precede all
double-primed ones, move the block $\eta'$, of degree $p'$, to the left past
the block $\xi$, of degree $q$.  Interchanging two homogeneous elements of
degrees $a$ and $b$ in an exterior algebra produces the sign $(-1)^{ab}$;
here this gives $(-1)^{qp'}$ and proves
(1), since $\eta\wedge\eta'$ and $\xi\wedge\xi'$
are the primed and double-primed blocks of the result.
Assertion (2) is the graded commutativity of
$\bigwedge^\bullet W$ applied to the homogeneous elements $\alpha$ and
$\beta$ of degrees $p+q$ and $p'+q'$; note that the bidegrees enter only
through their sum.  Assertion (3) holds because
the exterior algebra is associative and the identification of
\Cref{lem:ch09-exterior-of-double} is multiplicative; bilinearity over
$C^\infty(U)$ is the definition of the coefficientwise product.
\end{proof}

\begin{warning}{Total degree, not bidegree}
The commutation sign in \Cref{prop:ch09-wedge-signs}\,(2) is governed by the
\emph{total} degrees $p+q$ and $p'+q'$, not by the bidegrees separately.
Thus $(1,1)$-forms commute with everything of even total degree and, in
particular, with each other; this is what makes powers such as
$(\dprimeddprime u)^k$ well behaved.  By contrast the internal sign
$(-1)^{qp'}$ of (1) depends on the bidegrees, and
it is the sign that is forgotten in hand computations.  A reliable habit:
never move a labelled block past another without writing down its degree.
\end{warning}

For later sign checks, if $e_i=\dprime x_i$ and
$f_i=\ddprime x_i$, then
\[
 e_1\wedge f_1\wedge\cdots\wedge e_n\wedge f_n
 =(-1)^{n(n-1)/2}
 e_1\wedge\cdots\wedge e_n\wedge f_1\wedge\cdots\wedge f_n.
\]
The exponent counts the number of double-primed factors moved past later
primed factors.

\begin{lemma}
\label{lem:ch09-interlaced}
Write $c_k=(-1)^{k(k-1)/2}$ and, for a basis $x_1,\ldots,x_n$ of $V^*$,
\[
 \lambda=\dprime x_1\wedge\cdots\wedge\dprime x_n,
 \qquad
 \tau_n=\dprime x_1\wedge\ddprime x_1\wedge\cdots\wedge
        \dprime x_n\wedge\ddprime x_n .
\]
Then
\[
 \tau_n=c_n\,\dprime x_1\wedge\cdots\wedge\dprime x_n\wedge
              \ddprime x_1\wedge\cdots\wedge\ddprime x_n .
\]
Moreover $c_ac_b(-1)^{ab}=c_{a+b}$ for all $a,b\in\mathbf N$.
\end{lemma}

\begin{proof}
Put $e_i=\dprime x_i$ and $f_i=\ddprime x_i$ as above.  Starting from the
interlaced product $e_1f_1e_2f_2\cdots e_nf_n$, move $f_1$ to the right past
$e_2,\ldots,e_n$, which costs $(-1)^{n-1}$; then move $f_2$ past
$e_3,\ldots,e_n$, at a cost of $(-1)^{n-2}$; and so on, the last move being
free.  The total sign is
\[
 (-1)^{(n-1)+(n-2)+\cdots+1+0}=(-1)^{n(n-1)/2}=c_n,
\]
and after all moves the product is $e_1\cdots e_nf_1\cdots f_n$ with the
$f$'s in their original relative order.  Since $c_n^2=1$, the two displayed
forms of $\tau_n$ are equivalent.

For the numerical identity, compute the exponents modulo $2$:
\[
 \frac{a(a-1)}2+\frac{b(b-1)}2+ab
 =\frac{a^2-a+b^2-b+2ab}2
 =\frac{(a+b)^2-(a+b)}2 ,
\]
and the right-hand side is $\tfrac12(a+b)(a+b-1)$.
\end{proof}

The innocuous identity $c_ac_b(-1)^{ab}=c_{a+b}$ is used constantly below.
It is what makes the constants $c_p$ in the positivity generators fit
together under multiplication, and it is the reason that the complementary
pairing of \Cref{lem:ch09-complement-pairing} carries no sign at all.  With
$a=p$ and $b=n-p$ it reads
\[
 c_pc_{n-p}(-1)^{p(n-p)}=c_n .
\]

\begin{example}[Superforms in one variable]
\label{ex:ch09-line}
Let $V=\mathbf R$ with coordinate $x$ and let $U\subset\mathbf R$ be open.
There are exactly four spaces of superforms, of total dimension
$4=4^1$:
\[
 \begin{array}{llll}
 A^{0,0}(U)=C^\infty(U), &
 A^{1,0}(U)=C^\infty(U)\,\dprime x, \\[2pt]
 A^{0,1}(U)=C^\infty(U)\,\ddprime x, &
 A^{1,1}(U)=C^\infty(U)\,\dprime x\wedge\ddprime x .
 \end{array}
\]
All products of two generators vanish except
$\dprime x\wedge\ddprime x=-\ddprime x\wedge\dprime x$, and
$\tau_1=\dprime x\wedge\ddprime x$ with $c_1=1$.  For $f\in C^\infty(U)$,
\[
 \dprime f=f'\,\dprime x,\qquad
 \ddprime f=f'\,\ddprime x,\qquad
 \dprimeddprime f=f''\,\dprime x\wedge\ddprime x=f''\tau_1 .
\]
Thus in one variable the entire calculus reduces to the statement that
$\dprimeddprime f\geq0$ means $f''\geq0$.  The reader who keeps this case in
mind will not be misled by the signs in higher dimension: every formula below
must restrict to it.
\end{example}

\begin{example}[Superforms in two variables]
\label{ex:ch09-plane}
Let $V=\mathbf R^2$ with coordinates $x_1,x_2$.  The bigraded pieces have
dimensions $\binom2p\binom2q$, displayed as the matrix
\[
 \bigl(\dim A^{p,q}\bigr)_{0\le p,q\le2}
 =\begin{pmatrix}1&2&1\\2&4&2\\1&2&1\end{pmatrix},
 \qquad\text{total }16=4^2 .
\]
A general $(1,1)$-form is
\[
 \alpha=a_{11}\,\dprime x_1\wedge\ddprime x_1
       +a_{12}\,\dprime x_1\wedge\ddprime x_2
       +a_{21}\,\dprime x_2\wedge\ddprime x_1
       +a_{22}\,\dprime x_2\wedge\ddprime x_2
\]
with $a_{ij}\in C^\infty(U)$, and the top space $A^{2,2}(U)$ is spanned by
\[
 \tau_2=\dprime x_1\wedge\ddprime x_1\wedge\dprime x_2\wedge\ddprime x_2
      =-\,\dprime x_1\wedge\dprime x_2\wedge\ddprime x_1\wedge\ddprime x_2,
\]
the sign being $c_2=-1$ as in \Cref{lem:ch09-interlaced}.  Two products are
worth doing by hand.  First,
\[
 \begin{aligned}
  &\dprime x_1\wedge\ddprime x_2\wedge\dprime x_2\wedge\ddprime x_1\\
  &\qquad=-\,\dprime x_1\wedge\dprime x_2
             \wedge\ddprime x_2\wedge\ddprime x_1\\
  &\qquad=\dprime x_1\wedge\dprime x_2
             \wedge\ddprime x_1\wedge\ddprime x_2
   =-\tau_2 ,
 \end{aligned}
\]
where the first step moves $\ddprime x_2$ past $\dprime x_2$ and the second
transposes the two double-primed factors.  Second, and by the same two moves,
\[
 \dprime x_1\wedge\ddprime x_1\wedge\dprime x_2\wedge\ddprime x_2=\tau_2 .
\]
Hence for the $(1,1)$-form $\alpha$ above,
\[
 \alpha\wedge\alpha=2(a_{11}a_{22}-a_{12}a_{21})\,\tau_2
 =2\det(a_{ij})\,\tau_2 ,
\]
which is the case $n=2$ of \Cref{prop:superform-matrix-wedge}.  The factor
$2=2!$ is not a normalization one is free to change: it counts the two ways
of choosing which factor contributes the index $1$.
\end{example}

\subsection{The two differentials and the involution}

We now differentiate.  There is only one derivative available, the ordinary
one; what has to be decided is into which labelled block the resulting
one-form is inserted, and where in the product it is placed.

\begin{definition}
\label{def:ch09-operators}
For $\alpha\in A^{p,q}(U)$ define
\[
\begin{aligned}
\dprime\alpha&=\sum_{I,J,i}\frac{\partial\alpha_{IJ}}{\partial x_i}
 \dprime x_i\wedge \dprime x_I\wedge \ddprime x_J,\\
\ddprime\alpha&=(-1)^p\sum_{I,J,i}\frac{\partial\alpha_{IJ}}{\partial x_i}
 \dprime x_I\wedge \ddprime x_i\wedge \ddprime x_J.
\end{aligned}
\]
The involution $J:A^{p,q}\to A^{q,p}$ is the algebra automorphism determined
by $J(\dprime x_i)=\ddprime x_i$, $J(\ddprime x_i)=\dprime x_i$, and $Jf=f$.
\end{definition}

Three points about \Cref{def:ch09-operators} should be settled before it is
used.

The operator $\dprime$ raises the first degree and $\ddprime$ the second, so
$\dprime:A^{p,q}\to A^{p+1,q}$ and $\ddprime:A^{p,q}\to A^{p,q+1}$.  Both are
$\mathbf R$-linear but neither is $C^\infty(U)$-linear: they are first-order
differential operators, and their failure to be linear over functions is
precisely the Leibniz rule.

The sign $(-1)^p$ in the definition of $\ddprime$ is not decorative.  The new
factor $\ddprime x_i$ must be inserted at the head of the double-primed block,
which means moving it past the $p$ primed factors of $\dprime x_I$; the sign
$(-1)^p$ compensates for exactly that passage.  With this convention
$\ddprime$ is the antiderivation that one obtains from $\dprime$ by
conjugating with $J$, as \Cref{lem:ch09-J-conjugation} makes precise, and both
operators obey the same Leibniz rule with the same sign.  Without it one
would have to carry an extra $(-1)^p$ through every later formula.

Finally, the involution is well defined.  Indeed the linear map
$V^*\oplus V^*\to V^*\oplus V^*$ exchanging the two summands extends uniquely
to an automorphism of the exterior algebra $\bigwedge^\bullet(V^*\oplus V^*)$,
and \Cref{lem:ch09-exterior-of-double} transports that automorphism to
$A^{\bullet,\bullet}(U)$; extending it by the identity on coefficients gives
$J$.  Since the exchange map is an involution of $V^*\oplus V^*$, we get
\[
 J^2=\operatorname{id}\quad\text{on }A^{\bullet,\bullet}(U),
 \qquad
 J(\alpha\wedge\beta)=J\alpha\wedge J\beta .
\]
In coordinates, and with $|I|=p$, $|J|=q$,
\begin{equation}
\label{eq:ch09-J-coordinates}
 J\Bigl(\sum_{I,J}\alpha_{IJ}\,\dprime x_I\wedge\ddprime x_J\Bigr)
 =\sum_{I,J}\alpha_{IJ}\,\ddprime x_I\wedge\dprime x_J
 =(-1)^{pq}\sum_{I,J}\alpha_{IJ}\,\dprime x_J\wedge\ddprime x_I ,
\end{equation}
the last equality being the passage of the block $\ddprime x_I$, of degree
$p$, past the block $\dprime x_J$, of degree $q$.

On functions these conventions give
\[
 \dprime f=\sum_i f_i\,\dprime x_i,
 \qquad
 \ddprime f=\sum_i f_i\,\ddprime x_i,
 \qquad
 \dprimeddprime f=\sum_{i,j}f_{ij}\,
                 \dprime x_i\wedge\ddprime x_j.
\]
Thus $\dprimeddprime$ annihilates affine functions.  The order of the two
differentials matters: $\ddprime\dprime f=-\dprimeddprime f$.

\begin{proposition}
\label{prop:superform-differential-identities}
One has
\[
(\dprime)^2=(\ddprime)^2=0,\qquad \dprimeddprime+\ddprime\dprime=0,
\]
and the graded Leibniz rules
\[
\dprime(\alpha\wedge\beta)=\dprime\alpha\wedge\beta+
 (-1)^{p+q}\alpha\wedge \dprime\beta,
\]
with the analogous formula for $\ddprime$.
\end{proposition}
\begin{proof}
It is convenient to record first a reformulation of
\Cref{def:ch09-operators} that does not presuppose sorted multi-indices.
Call \emph{monomial} a product $\mu$ of generators $\dprime x_i$ and
$\ddprime x_j$, in any order, and let $m$ be its total degree.  Every
superform is a finite sum of terms $a\mu$ with $a\in C^\infty(U)$, and
\begin{equation}
\label{eq:ch09-dprime-monomial}
 \dprime(a\mu)=\sum_i\frac{\partial a}{\partial x_i}\,\dprime x_i\wedge\mu .
\end{equation}
This is consistent: re-sorting $\mu$ into the normal form
$\dprime x_I\wedge\ddprime x_J$ multiplies both sides of
\eqref{eq:ch09-dprime-monomial} by the same constant sign, and on normal-form
monomials \eqref{eq:ch09-dprime-monomial} is the first line of
\Cref{def:ch09-operators}.

\emph{Leibniz rule for $\dprime$.}  By bilinearity it suffices to take
$\alpha=a\mu$ and $\beta=b\nu$ with $\mu,\nu$ monomials, $\mu$ of total
degree $m=p+q$.  Then $\alpha\wedge\beta=ab\,\mu\wedge\nu$ and, by
\eqref{eq:ch09-dprime-monomial} and the product rule for partial derivatives,
\[
 \dprime(\alpha\wedge\beta)
 =\sum_i\Bigl(\frac{\partial a}{\partial x_i}b
       +a\frac{\partial b}{\partial x_i}\Bigr)
     \dprime x_i\wedge\mu\wedge\nu .
\]
The first group of terms is $\dprime\alpha\wedge\beta$.  In the second,
moving the degree-one factor $\dprime x_i$ past $\mu$ costs $(-1)^m$, so that
group equals
\[
 (-1)^{m}\,a\mu\wedge\sum_i\frac{\partial b}{\partial x_i}\dprime x_i\wedge\nu
 =(-1)^{p+q}\alpha\wedge\dprime\beta .
\]
This is the asserted rule.

\emph{Vanishing of $(\dprime)^2$.}  With $\alpha=a\mu$,
\[
 \dprime\dprime\alpha
 =\sum_{i,j}\frac{\partial^2a}{\partial x_j\partial x_i}\,
   \dprime x_j\wedge\dprime x_i\wedge\mu .
\]
The coefficients are symmetric in $(i,j)$ by the equality of mixed partial
derivatives for $C^\infty$ functions, whereas
$\dprime x_j\wedge\dprime x_i=-\dprime x_i\wedge\dprime x_j$; the terms
therefore cancel in pairs and the diagonal terms vanish.

\emph{The operators $\ddprime$ and $J\dprime J$.}  Let
$\alpha=a\,\dprime x_I\wedge\ddprime x_J$ with $|I|=p$ and $|J|=q$.  By
\eqref{eq:ch09-J-coordinates}, $J\alpha=(-1)^{pq}a\,\dprime x_J\wedge\ddprime x_I$,
whence
\[
 \dprime J\alpha
 =(-1)^{pq}\sum_i\frac{\partial a}{\partial x_i}\,
   \dprime x_i\wedge\dprime x_J\wedge\ddprime x_I .
\]
The monomial appearing here has a primed block of degree $q+1$ and a
double-primed block of degree $p$, so applying
\eqref{eq:ch09-J-coordinates} once more produces the sign $(-1)^{(q+1)p}$ and
gives
\[
\begin{aligned}
 J\dprime J\alpha
 &=(-1)^{pq+(q+1)p}\sum_i\frac{\partial a}{\partial x_i}\,
   \dprime x_I\wedge\ddprime x_i\wedge\ddprime x_J\\
 &=(-1)^{p}\sum_i\frac{\partial a}{\partial x_i}\,
   \dprime x_I\wedge\ddprime x_i\wedge\ddprime x_J
 =\ddprime\alpha,
\end{aligned}
\]
because $(-1)^{2pq}=1$.  Consequently
$\ddprime=J\dprime J$, and since $J$ is an algebra automorphism preserving
total degrees, the Leibniz rule for $\ddprime$ follows from that for
$\dprime$:
\[
\begin{aligned}
 \ddprime(\alpha\wedge\beta)
 &=J\dprime\bigl(J\alpha\wedge J\beta\bigr)
 =J\bigl(\dprime J\alpha\wedge J\beta
   +(-1)^{p+q}J\alpha\wedge\dprime J\beta\bigr)\\
 &=\ddprime\alpha\wedge\beta+(-1)^{p+q}\alpha\wedge\ddprime\beta .
\end{aligned}
\]
Likewise $(\ddprime)^2=J(\dprime)^2J=0$.

\emph{Anticommutation.}  Again with $\alpha=a\,\dprime x_I\wedge\ddprime x_J$
of bidegree $(p,q)$,
\[
 \dprimeddprime\alpha
 =(-1)^{p}\sum_{i,j}\frac{\partial^2a}{\partial x_i\partial x_j}\,
   \dprime x_i\wedge\dprime x_I\wedge\ddprime x_j\wedge\ddprime x_J,
\]
while $\dprime\alpha$ has bidegree $(p+1,q)$, so that
\[
 \ddprime\dprime\alpha
 =(-1)^{p+1}\sum_{i,j}\frac{\partial^2a}{\partial x_j\partial x_i}\,
   \dprime x_i\wedge\dprime x_I\wedge\ddprime x_j\wedge\ddprime x_J .
\]
The two sums have the same monomials and, by the equality of mixed partial
derivatives, the same coefficients; only the prefactors $(-1)^p$ and
$(-1)^{p+1}$ differ.  Hence $\dprimeddprime\alpha+\ddprime\dprime\alpha=0$.
\end{proof}

Two features of this proof are worth isolating, because they explain the
shape of the convention.  The vanishing of $(\dprime)^2$ is a genuine
analytic fact: it uses the symmetry of second derivatives.  The
anticommutation $\dprimeddprime=-\ddprime\dprime$, by contrast, is almost
entirely a matter of the sign $(-1)^p$ built into $\ddprime$; had that sign
been omitted, the two mixed compositions would have been equal instead of
opposite, and $\dprimeddprime$ would not be an antiderivation-like operator of
the type needed for Stokes' formula in
\Cref{chap:tropical-integration}.

\begin{lemma}
\label{lem:ch09-J-conjugation}
On $A^{p,q}(U)$ one has
\[
 J\dprime J=\ddprime,\qquad
 J\ddprime J=\dprime,\qquad
 J\dprime=\ddprime J,\qquad
 J\ddprime=\dprime J,
\]
and consequently $J\dprimeddprime J=-\dprimeddprime$.  If moreover
$\alpha\in A^{p,p}(U)$ satisfies $J\alpha=(-1)^p\alpha$, then
\[
 J(\dprime\alpha)=(-1)^p\ddprime\alpha,
 \qquad
 J(\ddprime\alpha)=(-1)^p\dprime\alpha .
\]
\end{lemma}

\begin{proof}
The identity $J\dprime J=\ddprime$ was proved inside
\Cref{prop:superform-differential-identities}.  Composing it on both sides
with $J$ and using $J^2=\operatorname{id}$ gives $\dprime=J\ddprime J$ and the
two commutation forms.  For the third assertion,
\[
 J\dprimeddprime J
 =(J\dprime J)(J\ddprime J)
 =\ddprime\dprime
 =-\dprimeddprime ,
\]
where the first equality inserts $J^2=\operatorname{id}$ between the two
operators and the last is
\Cref{prop:superform-differential-identities}.  Finally, if
$J\alpha=(-1)^p\alpha$ then
$J\dprime\alpha=\ddprime J\alpha=(-1)^p\ddprime\alpha$, and symmetrically.
\end{proof}

\begin{corollary}
\label{cor:ch09-symmetric-exchange}
Let $p+q+r=n-1$ and let $\alpha,\beta,\gamma$ be superforms of bidegrees
$(p,p)$, $(q,q)$, $(r,r)$ with $J\alpha=(-1)^p\alpha$,
$J\beta=(-1)^q\beta$, $J\gamma=(-1)^r\gamma$.  Then
\[
 \dprime\alpha\wedge\ddprime\beta\wedge\gamma
 =-\,\ddprime\alpha\wedge\dprime\beta\wedge\gamma .
\]
\end{corollary}

\begin{proof}
Both sides are superforms of bidegree
$(p+q+r+1,p+q+r+1)=(n,n)$.  The space $A^{n,n}(U)$ is free of rank one over
$C^\infty(U)$, generated by $\tau_n$, and
\eqref{eq:ch09-J-coordinates} with $I=J=\{1,\ldots,n\}$ gives
$J\tau_n=(-1)^n\tau_n$; hence every $(n,n)$-form $\omega$ satisfies
$J\omega=(-1)^n\omega$.  Applying $J$ to the left-hand side and using that
$J$ is an algebra automorphism together with
\Cref{lem:ch09-J-conjugation},
\[
\begin{aligned}
 J\bigl(\dprime\alpha\wedge\ddprime\beta\wedge\gamma\bigr)
 &=J\dprime\alpha\wedge J\ddprime\beta\wedge J\gamma
 =(-1)^{p}(-1)^{q}(-1)^{r}\,
   \ddprime\alpha\wedge\dprime\beta\wedge\gamma\\
 &=(-1)^{n-1}\ddprime\alpha\wedge\dprime\beta\wedge\gamma .
\end{aligned}
\]
Comparing with $J\omega=(-1)^n\omega$ for
$\omega=\dprime\alpha\wedge\ddprime\beta\wedge\gamma$ yields
$(-1)^n\omega=(-1)^{n-1}\ddprime\alpha\wedge\dprime\beta\wedge\gamma$, which
is the assertion.
\end{proof}

\Cref{cor:ch09-symmetric-exchange} is the exchange rule behind the Green
formulae of \Cref{chap:tropical-integration}: it is what allows one to move a
prime from one factor to the other at the cost of a single sign, and it is
the reason that symmetric forms, rather than arbitrary ones, are the natural
inputs of the integration theory.

\begin{remark}[Two involutions]
\label{rem:ch09-two-involutions}
The revised Chambert--Loir--Ducros manuscript \cite[\S1.1]{CLD25} uses a
different involution, namely the algebra automorphism $J_{\mathrm{CLD}}$
determined by
\[
 J_{\mathrm{CLD}}(\dprime x)=\ddprime x,
 \qquad
 J_{\mathrm{CLD}}(\ddprime x)=-\dprime x
 \qquad(x\in V^*),
\]
which is the doubled analogue of a complex structure rather than an
involution: one has $J_{\mathrm{CLD}}^2=(-1)^{p+q}$ on $A^{p,q}$.  The
operator used in this book is Lagerberg's original one, with
$J(\ddprime x)=\dprime x$ and $J^2=\operatorname{id}$.  The two are related
by
\[
 J_{\mathrm{CLD}}=(-1)^qJ\quad\text{on }A^{p,q},
\]
as one checks on generators, and this rescaling is multiplicative because
$(-1)^{q}(-1)^{q'}=(-1)^{q+q'}$.  Three consequences deserve emphasis.
\begin{enumerate}
\item The two conventions define the \emph{same} notion of symmetry: on
  $A^{p,p}$ the condition $J_{\mathrm{CLD}}\omega=\omega$ is equivalent to
  $J\omega=(-1)^p\omega$, and both say that the coefficient array satisfies
  $\omega_{IJ}=\omega_{JI}$ (\Cref{prop:ch09-symmetry-matrix}).
\item The two conventions agree on $A^{p,0}$, since neither introduces a sign
  when a primed factor is converted into a double-primed one.  The generators
  $\eta\wedge J\eta$ of the positivity cones are therefore literally the same
  in both, and \emph{no positivity statement in this book depends on the
  choice}.
\item The commutation rules differ in shape but not in content:
  \Cref{lem:ch09-J-conjugation} gives $J\dprime=\ddprime J$ and
  $J\ddprime=\dprime J$, while the same computation with
  $J_{\mathrm{CLD}}$ gives
  $J_{\mathrm{CLD}}\dprime=\ddprime J_{\mathrm{CLD}}$ and
  $J_{\mathrm{CLD}}\ddprime=-\dprime J_{\mathrm{CLD}}$, whence
  $J_{\mathrm{CLD}}\dprimeddprime=\dprimeddprime J_{\mathrm{CLD}}$ where we
  have $J\dprimeddprime J=-\dprimeddprime$.
\end{enumerate}
The reader consulting \cite{CLD25} alongside this chapter should translate by
the displayed rule and should expect isolated signs, never a different
theorem.
\end{remark}

\subsection{The mixed operator on functions}

The composite $\dprimeddprime$ is the operator of interest in potential
theory, and its behaviour on functions is elementary but must be recorded
with care, since the two Leibniz rules interact.

\begin{lemma}
\label{lem:ch09-ddprime-product}
For $u,v\in C^\infty(U)$,
\[
 \dprimeddprime(uv)
 =u\,\dprimeddprime v+v\,\dprimeddprime u
  +\dprime u\wedge\ddprime v+\dprime v\wedge\ddprime u .
\]
In particular $\dprimeddprime$ is not a derivation, and the two mixed terms
do not combine: $\dprime u\wedge\ddprime v$ and $\dprime v\wedge\ddprime u$
are different superforms unless $u$ and $v$ have proportional differentials.
\end{lemma}

\begin{proof}
By the Leibniz rule of \Cref{prop:superform-differential-identities} in
bidegree $(0,0)$,
\[
 \ddprime(uv)=u\,\ddprime v+v\,\ddprime u .
\]
Apply $\dprime$ and use the Leibniz rule once more, now with a factor of
total degree zero in front of a factor of total degree one:
\[
 \dprime\bigl(u\,\ddprime v\bigr)
 =\dprime u\wedge\ddprime v+u\,\dprimeddprime v,
\]
and symmetrically for $v\,\ddprime u$.  Adding the two identities gives the
formula.  For the last sentence, observe that in coordinates the coefficient
matrix of $\dprime u\wedge\ddprime v$ is
$(\partial_iu)(\partial_jv)$, whose transpose is the coefficient matrix of
$\dprime v\wedge\ddprime u$; the two agree exactly when the rank-one matrix
$\nabla u\,\nabla v^{\mathsf T}$ is symmetric.
\end{proof}

\begin{lemma}
\label{lem:superform-composition}
Let $f\in C^\infty(U)$ and let $\chi$ be smooth on an interval containing
$f(U)$. Then
\[
 \dprimeddprime(\chi\circ f)
 =\chi'(f)\,\dprimeddprime f
  +\chi''(f)\,\dprime f\wedge\ddprime f.
\]
\end{lemma}
\begin{proof}
Apply $\dprime$ to
$\ddprime(\chi\circ f)=\chi'(f)\ddprime f$ and use the Leibniz rule.  No
extra sign occurs because $\chi'(f)$ has total degree zero.

In detail: the chain rule for the ordinary derivative gives
$\partial_i(\chi\circ f)=\chi'(f)\partial_if$, that is,
$\ddprime(\chi\circ f)=\chi'(f)\,\ddprime f$.  Differentiating again with
$\dprime$ and using \Cref{prop:superform-differential-identities} with
$\alpha=\chi'(f)$ of bidegree $(0,0)$,
\[
 \dprimeddprime(\chi\circ f)
 =\dprime\bigl(\chi'(f)\bigr)\wedge\ddprime f
   +\chi'(f)\,\dprimeddprime f
 =\chi''(f)\,\dprime f\wedge\ddprime f+\chi'(f)\,\dprimeddprime f,
\]
where the last step applies the chain rule to $\chi'\circ f$.  The sign
$(-1)^{p+q}$ of the Leibniz rule is $+1$ here because $p=q=0$.
\end{proof}

The two terms in \Cref{lem:superform-composition} have a transparent meaning
in coordinates: the coefficient matrix of the left-hand side is
\[
 \chi'(f)\,\bigl(\partial_i\partial_jf\bigr)
 +\chi''(f)\,\bigl(\partial_if\bigr)\bigl(\partial_jf\bigr),
\]
which is the classical formula for the Hessian of a composition.  The
superform calculus adds nothing to that computation; what it adds is the
observation that the second summand is a rank-one form of a special shape,
namely $\eta\wedge J\eta$ with $\eta=\dprime f$, and such forms are exactly
the generators of the smallest positivity cone.  This is the mechanism behind
\Cref{cor:superform-convex-composition}.

\begin{example}
\label{ex:ch09-hessian-form}
For $f\in C^\infty(U)$,
\[
\dprimeddprime f=\sum_{i,j}\frac{\partial^2f}{\partial x_i\partial x_j}
                      \dprime x_i\wedge \ddprime x_j.
\]
On $\mathbf R$, $\dprimeddprime f=f''\,\dprime x\wedge \ddprime x$. This formula is our first sign
test and will later identify convex functions with psh functions.
\end{example}

This is the normalization recorded as \eqref{eq:appJ-real-hessian} in
\Cref{sec:appJ-cld}, and it is binding for the whole book: there is no factor
$\tfrac12$, no $2\pi$, and no imaginary unit.  The reader coming from the
complex side should note that the corresponding complex operator carries the
normalization $\ddc=\tfrac{\rmi}{2\pi}\partial\bar\partial$ fixed in the
conventions; the comparison of the two, and the resulting factors in the
Poincar\'e--Lelong formulae, is made once and for all in
\Cref{app:normalizations} and is never performed silently.

\begin{example}[A two-variable sign computation]
\label{ex:ch09-two-variable-sign}
For $f(x,y)=x^2+xy+2y^2$ one has
\[
 \dprimeddprime f
 =2\,\dprime x\wedge\ddprime x
  +\dprime x\wedge\ddprime y
  +\dprime y\wedge\ddprime x
  +4\,\dprime y\wedge\ddprime y.
\]
Its coefficient matrix has determinant $7$.  Expanding the wedge square with
the interlaced orientation gives
\[
 (\dprimeddprime f)^2
 =14\,\dprime x\wedge\ddprime x\wedge
       \dprime y\wedge\ddprime y.
\]
This is the first instance of the determinant identity proved below.
\end{example}

It is instructive to expand the last computation by hand, since it exhibits
every sign of the chapter in miniature.  Write $\alpha=\dprimeddprime f$ and
abbreviate $e_1=\dprime x$, $e_2=\dprime y$, $f_1=\ddprime x$,
$f_2=\ddprime y$, so that
\[
 \alpha=2e_1f_1+e_1f_2+e_2f_1+4e_2f_2 .
\]
Each summand has total degree two, so all summands commute with one another
and $\alpha\wedge\alpha$ may be expanded as if the factors were scalars.  A
product $e_if_j\wedge e_kf_l$ vanishes unless $\{i,k\}=\{j,l\}=\{1,2\}$, and
by \Cref{ex:ch09-plane} we have $e_1f_1\wedge e_2f_2=\tau_2$ and
$e_1f_2\wedge e_2f_1=-\tau_2$.  Hence
\[
 \alpha\wedge\alpha
 =2\bigl(2\cdot4\bigr)\tau_2+2\bigl(1\cdot1\bigr)(-\tau_2)
 =\bigl(16-2\bigr)\tau_2=14\,\tau_2,
\]
the outer factors $2$ counting the two orders in which the two chosen
summands can be multiplied.  This is $2!\det\binom{2\ 1}{1\ 4}=2\cdot7$, as
asserted.  The single minus sign in $e_1f_2\wedge e_2f_1=-\tau_2$ is the
entire content of the determinant: it converts a naive product of diagonal
terms into an alternating sum.

\section{Contraction in the two labelled blocks}
\label{sec:ch09-contraction}

Differentiation raises degrees; contraction lowers them.  Since a superform
has two blocks, there are two contraction operators, and---unlike
$\dprime$ and $\ddprime$---the one acting on the second block is taken here
\emph{without} the sign that would record the passage across the first.  That
choice is a convention: it is the one in which the boundary formulas of
\Cref{chap:tropical-integration} are written, the alternative differs from it
by $(-1)^p$, and both are set out below.  These operators are not needed
for the algebra of forms, but every boundary formula in
\Cref{chap:tropical-integration} is written in them, as are the homogeneity
conditions \eqref{eq:appD-strong-homogeneity} for the support currents of
convex bodies, and their signs are exactly of the kind that this chapter
exists to pin down.

Recall that for $v\in V$ the interior product
$\iota_v:\bigwedge^kV^*\to\bigwedge^{k-1}V^*$ is the unique antiderivation of
degree $-1$ with $\iota_v\xi=\xi(v)$ for $\xi\in V^*$.

\begin{definition}
\label{def:ch09-contraction}
For $v\in V$ define the \emph{block contractions}
$\iota'_v:A^{p,q}(U)\to A^{p-1,q}(U)$ and
$\iota''_v:A^{p,q}(U)\to A^{p,q-1}(U)$ by
\[
 \begin{aligned}
  \iota'_v\bigl(a\,\dprime x_I\wedge\ddprime x_J\bigr)
   &=a\,(\iota_v\dprime x_I)\wedge\ddprime x_J,\\
  \iota''_v\bigl(a\,\dprime x_I\wedge\ddprime x_J\bigr)
   &=a\,\dprime x_I\wedge(\iota_v\ddprime x_J),
 \end{aligned}
\]
with $|I|=p$, extended $C^\infty(U)$-linearly.  Here $\iota_v$ acts on the
primed, respectively double-primed, exterior factor through the
identification of each with $\bigwedge^\bullet V^*$; neither operator carries
a traversal sign, each simply ignoring the block on which it does not act.
Alongside $\iota''_v$ we also use the \emph{graded contraction}
\[
 \widetilde\iota''_v:=(-1)^p\,\iota''_v
 \qquad\text{on }A^{p,q}(U),
\]
which, unlike $\iota''_v$, is an interior product for the total degree.
\end{definition}

\begin{remark}
\label{rem:ch09-contraction-convention}
The two normalizations of the double-primed contraction differ by the
traversal sign $(-1)^p$, and both occur in the literature.  This book uses
the block convention throughout: it is the convention in which the boundary
traces of \Cref{chap:tropical-integration} are stated
(\Cref{def:ch11-contractions} and \eqref{eq:superform-boundary-traces}), and
it is recorded in \Cref{sec:appJ-currents}.  The graded variant
$\widetilde\iota''_v$ is kept in play because it, and not $\iota''_v$, obeys
the clean structural identities of \Cref{lem:ch09-contraction-basic} and
\Cref{prop:ch09-cartan}.  Wherever a factor $(-1)^p$ appears beside
$\iota''_v$ below, it is the conversion between the two.
\end{remark}

\begin{lemma}
\label{lem:ch09-contraction-basic}
Let $v,w\in V$.
\begin{enumerate}
\item Under the isomorphism of \Cref{lem:ch09-exterior-of-double}, $\iota'_v$
  and $\widetilde\iota''_v$ are the interior products of
  $\bigwedge^\bullet(V^*\oplus V^*)$ by the vectors $(v,0)$ and $(0,v)$ of
  $V\oplus V$.  In particular both are antiderivations for the total degree:
  \[
   \iota'_v(\alpha\wedge\beta)
   =\iota'_v\alpha\wedge\beta+(-1)^{p+q}\alpha\wedge\iota'_v\beta,
   \qquad\alpha\in A^{p,q}(U),
  \]
  and likewise for $\widetilde\iota''_v$.  The block operator $\iota''_v$ is
  an antiderivation within its own block only; its asymmetric Leibniz rule is
  recorded in \Cref{lem:ch11-contraction-basic}.
\item $\iota'_v\iota'_w=-\iota'_w\iota'_v$ and
  $\iota''_v\iota''_w=-\iota''_w\iota''_v$, so
  $(\iota'_v)^2=(\iota''_v)^2=0$.  The mixed pairs \emph{commute},
  \[
   \iota'_v\iota''_w=\iota''_w\iota'_v ,
   \qquad\text{equivalently}\qquad
   \iota'_v\widetilde\iota''_w=-\widetilde\iota''_w\iota'_v ,
  \]
  the first because the two operators act on different blocks.
\item On $A^{p,q}(U)$,
  \[
   J\iota'_vJ=(-1)^p\iota''_v=\widetilde\iota''_v,
   \qquad
   J\iota''_vJ=(-1)^q\iota'_v .
  \]
\end{enumerate}
\end{lemma}

\begin{proof}
(1)  The interior product of $\bigwedge^\bullet(V^*\oplus V^*)$ by
$(v,0)$ annihilates every double-primed generator and sends
$\dprime x$ to $x(v)$.  Being an antiderivation, it therefore satisfies
\[
 \iota_{(v,0)}\bigl(\dprime x_I\wedge\ddprime x_J\bigr)
 =(\iota_v\dprime x_I)\wedge\ddprime x_J
  +(-1)^{|I|}\dprime x_I\wedge\iota_{(v,0)}\ddprime x_J
 =(\iota_v\dprime x_I)\wedge\ddprime x_J ,
\]
which is $\iota'_v$.  Similarly $\iota_{(0,v)}$ annihilates the primed
generators, so
\[
 \iota_{(0,v)}\bigl(\dprime x_I\wedge\ddprime x_J\bigr)
 =(-1)^{|I|}\dprime x_I\wedge(\iota_v\ddprime x_J)
 =(-1)^{p}\,\iota''_v\bigl(\dprime x_I\wedge\ddprime x_J\bigr),
\]
which is $\widetilde\iota''_v$.  The sign $(-1)^p$ separating the two
operators is exactly the antiderivation sign incurred in crossing the primed
block; the block operator $\iota''_v$ discards it by fiat, which is precisely
why $\iota''_v$ is not an antiderivation for the total degree.

(2) For $\iota'_v$ and $\widetilde\iota''_v$ all three relations, including
$\iota'_v\widetilde\iota''_w=-\widetilde\iota''_w\iota'_v$, are the standard
identity $\iota_a\iota_b=-\iota_b\iota_a$ for interior products, applied to
$a,b\in V\oplus V$.  Passing to the block operator, $\iota''_v$ preserves the
primed degree, so the two conversion factors $(-1)^p$ in
$\widetilde\iota''_v\widetilde\iota''_w=-\widetilde\iota''_w
\widetilde\iota''_v$ agree on the two sides and cancel, leaving
$\iota''_v\iota''_w=-\iota''_w\iota''_v$.  In the mixed relation, by
contrast, $\iota'_v$ lowers the primed degree by one, so the conversions on
the two sides are $(-1)^p$ and $(-1)^{p-1}$; the resulting extra minus sign
absorbs the one displayed, and $\iota'_v\iota''_w=\iota''_w\iota'_v$.

(3) Let $\alpha=a\,\dprime x_I\wedge\ddprime x_J$ with $|I|=p$, $|J|=q$.
By \eqref{eq:ch09-J-coordinates},
$J\alpha=(-1)^{pq}a\,\dprime x_J\wedge\ddprime x_I$, so
\[
 \iota'_vJ\alpha
 =(-1)^{pq}a\,(\iota_v\dprime x_J)\wedge\ddprime x_I ,
\]
a form whose primed block has degree $q-1$ and whose double-primed block has
degree $p$.  Applying $J$ and using \eqref{eq:ch09-J-coordinates} again
introduces the sign $(-1)^{(q-1)p}$, so
\[
 \begin{aligned}
 J\iota'_vJ\alpha
  &=(-1)^{pq+(q-1)p}a\,\dprime x_I\wedge(\iota_v\ddprime x_J)\\
  &=(-1)^{p}a\,\dprime x_I\wedge(\iota_v\ddprime x_J)\\
  &=(-1)^p\iota''_v\alpha
   =\widetilde\iota''_v\alpha .
 \end{aligned}
\]
The second identity is the same computation with the roles of the two blocks
interchanged, which replaces $p$ by $q$.
\end{proof}

The interaction with the differentials is the expected Cartan formula, with
the mixed pairs anticommuting.  For $v\in V$ write $\partial_v$ for the
directional derivative acting on coefficients, that is,
$\partial_v(a\mu)=(\partial_va)\mu$ for a monomial $\mu$; this is well
defined because the generators are constant.

\begin{proposition}
\label{prop:ch09-cartan}
For every $v\in V$ the graded contractions satisfy
\[
 \iota'_v\dprime+\dprime\iota'_v=\partial_v,
 \qquad
 \widetilde\iota''_v\ddprime+\ddprime\widetilde\iota''_v=\partial_v,
\]
\[
 \iota'_v\ddprime+\ddprime\iota'_v=0,
 \qquad
 \widetilde\iota''_v\dprime+\dprime\widetilde\iota''_v=0 .
\]
Equivalently, in terms of the block contraction and on $A^{p,q}(U)$,
\[
 \iota''_v\ddprime+\ddprime\iota''_v=(-1)^p\partial_v,
 \qquad
 \iota''_v\dprime-\dprime\iota''_v=0 .
\]
\end{proposition}

\begin{proof}
Let $\alpha=a\mu$ with $\mu$ a monomial of total degree $m$, and write
$v=\sum_iv_i\partial/\partial x_i$ in the coordinates at hand.  Using
\eqref{eq:ch09-dprime-monomial} and the antiderivation property of
$\iota'_v$,
\[
 \iota'_v\dprime\alpha
 =\sum_i\frac{\partial a}{\partial x_i}\,
   \iota'_v\bigl(\dprime x_i\wedge\mu\bigr)
 =\sum_i\frac{\partial a}{\partial x_i}
   \Bigl(v_i\,\mu-\dprime x_i\wedge\iota'_v\mu\Bigr)
 =\partial_va\cdot\mu
   -\sum_i\frac{\partial a}{\partial x_i}\dprime x_i\wedge\iota'_v\mu .
\]
On the other hand $\iota'_v\alpha=a\,\iota'_v\mu$, so
\[
 \dprime\iota'_v\alpha
 =\sum_i\frac{\partial a}{\partial x_i}\,\dprime x_i\wedge\iota'_v\mu .
\]
Adding gives $\partial_v\alpha$, which is the first identity; conjugating by
$J$ and using \Cref{lem:ch09-J-conjugation} together with
\Cref{lem:ch09-contraction-basic}(3) gives the second, since $J$ commutes
with $\partial_v$.

For the third, take $\alpha=a\,\dprime x_I\wedge\ddprime x_J$ with $|I|=p$.
Then
\[
 \iota'_v\ddprime\alpha
 =(-1)^{p}\sum_j\frac{\partial a}{\partial x_j}\,
   (\iota_v\dprime x_I)\wedge\ddprime x_j\wedge\ddprime x_J ,
\]
whereas $\iota'_v\alpha$ has first degree $p-1$, so
\[
 \ddprime\iota'_v\alpha
 =(-1)^{p-1}\sum_j\frac{\partial a}{\partial x_j}\,
   (\iota_v\dprime x_I)\wedge\ddprime x_j\wedge\ddprime x_J .
\]
These cancel.  The fourth identity is obtained from the third by conjugating
with $J$.

For the block forms, substitute $\widetilde\iota''_v=(-1)^p\iota''_v$ and
note that $\ddprime$ leaves the primed degree unchanged while $\dprime$
raises it by one.  In the second identity the two conversion factors are both
$(-1)^p$ and factor out, giving
$\iota''_v\ddprime+\ddprime\iota''_v=(-1)^p\partial_v$; in the fourth they are
$(-1)^{p+1}$ and $(-1)^p$, so the sum becomes a difference and
$\iota''_v\dprime=\dprime\iota''_v$.
\end{proof}

\begin{example}[Contractions of the Hessian form and of the top form]
\label{ex:ch09-contraction-examples}
Let $u\in C^\infty(U)$ and $v\in V$.  Contracting
$\dprimeddprime u=\sum_{i,j}u_{ij}\dprime x_i\wedge\ddprime x_j$ gives
\[
 \begin{aligned}
  \iota'_v\bigl(\dprimeddprime u\bigr)
   &=\sum_{i,j}u_{ij}v_i\,\ddprime x_j
     =\ddprime(\partial_vu),\\
    \iota''_v\bigl(\dprimeddprime u\bigr)
   &=\sum_{i,j}u_{ij}v_j\,\dprime x_i
     =\dprime(\partial_vu),
 \end{aligned}
\]
with no sign in the second formula, the block convention of
\Cref{def:ch09-contraction} leaving the primed factor untouched; the graded
contraction returns the opposite sign,
$\widetilde\iota''_v(\dprimeddprime u)=-\dprime(\partial_vu)$, since $p=1$
here.  Both identities also follow from \Cref{prop:ch09-cartan}, since
$\iota'_v\ddprime u=0$ and $\iota''_v\dprime u=0$.

For the top form, let $n=2$ and $v=\partial/\partial x_2$.  Regroup
\[
 \tau_2=\dprime x_1\wedge\ddprime x_1\wedge\dprime x_2\wedge\ddprime x_2
 =-\,\dprime x_1\wedge\dprime x_2\wedge\ddprime x_1\wedge\ddprime x_2
\]
and contract in the double-primed block alone, where
$\iota_{\partial/\partial x_2}(\ddprime x_1\wedge\ddprime x_2)
=-\,\ddprime x_1$.  This gives
\[
 \iota''_{\partial/\partial x_2}\tau_2
 =\dprime x_1\wedge\dprime x_2\wedge\ddprime x_1
 =-\,\dprime x_1\wedge\ddprime x_1\wedge\dprime x_2 ,
\]
and since $p=2$ the graded contraction returns the same form.  The boundary
trace \eqref{eq:superform-boundary-traces} of
\Cref{chap:tropical-integration} contracts the same block---the choice is
forced by the bidegree, a form of bidegree $(n-1,n)$ having no spare primed
factor---but carries in addition the prefactor $(-1)^{n-1}$, which comes from
the interlacing constant $c_nc_{n-1}$ and not from the contraction.
\end{example}

\section{Pullback and coordinate independence}
\label{sec:ch09-pullback}

Let $F:U\to U'$ be the restriction of an affine map. Define $F^*$ on
functions ordinarily and set
\[
F^*(\dprime y_i)=\dprime(y_i\circ F),\qquad
F^*(\ddprime y_i)=\ddprime(y_i\circ F),
\]
then extend multiplicatively.

More intrinsically: if $F:V\to V'$ is affine with linear part
$\overrightarrow F$, then the transpose
$\overrightarrow F{}^*:V'^*\to V^*$ induces a map
$\bigwedge^p V'^*\to\bigwedge^pV^*$ in each block, and
\[
 F^*\bigl(a\otimes\eta\otimes\xi\bigr)
 =(a\circ F)\otimes
   \bigl(\textstyle\bigwedge^p\overrightarrow F{}^*\eta\bigr)\otimes
   \bigl(\textstyle\bigwedge^q\overrightarrow F{}^*\xi\bigr),
 \qquad a\in C^\infty(U'),
\]
on $F^{-1}(U')$.  Under \Cref{lem:ch09-exterior-of-double} this is the map
induced on $\bigwedge^\bullet$ by
$\overrightarrow F{}^*\oplus\overrightarrow F{}^*:V'^*\oplus V'^*\to
V^*\oplus V^*$, so it is well defined and multiplicative before any
coordinate is chosen.  In particular $F^*$ does not depend on the coordinate
description, and it preserves bidegrees --- a point on which the whole
formalism turns, since an arbitrary smooth map would not preserve them in any
useful sense.

Write $F(x)=Bx+b$ in affine coordinates.  If
\[
 \alpha=\sum_{k,l}a_{kl}(y)\,
               \dprime y_k\wedge\ddprime y_l,
\]
then the coefficient matrix of $F^*\alpha$ at $x$ is
\[
                         B^{\mathsf T}A(F(x))B.
\]
This single formula simultaneously explains preservation of symmetry and of
positivity in bidegree $(1,1)$.  When $f$ is smooth on $U'$, it also yields
\[
 \operatorname{Hess}(f\circ F)
 =B^{\mathsf T}(\operatorname{Hess}f)\circ F\,B.
\]

\begin{lemma}
\label{lem:ch09-pullback-eleven}
Let $F(x)=Bx+b$ map $U\subset V=\mathbf R^n$ into
$U'\subset V'=\mathbf R^m$, where $B$ is the $m\times n$ matrix of the linear
part in the chosen coordinates.  If
$\alpha=\sum_{k,l}a_{kl}\,\dprime y_k\wedge\ddprime y_l$ has coefficient
matrix $A$, then
\[
 F^*\alpha
 =\sum_{i,j}\bigl(B^{\mathsf T}A(F(x))B\bigr)_{ij}\,
   \dprime x_i\wedge\ddprime x_j .
\]
Consequently, if $A$ is symmetric, so is $B^{\mathsf T}AB$; if $A$ is
positive semidefinite, so is $B^{\mathsf T}AB$; and for $f\in C^2(U')$,
\[
 \operatorname{Hess}(f\circ F)(x)
 =B^{\mathsf T}\operatorname{Hess}f(F(x))\,B .
\]
\end{lemma}

\begin{proof}
Since $y_k\circ F=\sum_iB_{ki}x_i+b_k$ and the differentials annihilate
constants, $F^*(\dprime y_k)=\sum_iB_{ki}\dprime x_i$ and
$F^*(\ddprime y_l)=\sum_jB_{lj}\ddprime x_j$.  Multiplicativity of $F^*$ and
bilinearity of the wedge product give
\[
 \begin{aligned}
  F^*\alpha
   &=\sum_{k,l}a_{kl}(F(x))\sum_{i,j}B_{ki}B_{lj}\,
     \dprime x_i\wedge\ddprime x_j\\
   &=\sum_{i,j}\Bigl(\sum_{k,l}B_{ki}a_{kl}(F(x))B_{lj}\Bigr)
     \dprime x_i\wedge\ddprime x_j,
 \end{aligned}
\]
and the inner sum is the $(i,j)$ entry of $B^{\mathsf T}A(F(x))B$.  The
matrix $B^{\mathsf T}AB$ is symmetric when $A$ is, and satisfies
$w^{\mathsf T}B^{\mathsf T}ABw=(Bw)^{\mathsf T}A(Bw)\geq0$ when $A$ is
positive semidefinite.  The Hessian statement is the case
$\alpha=\dprimeddprime f$, together with
\Cref{prop:ch09-pullback-functorial} below, which identifies
$F^*\dprimeddprime f$ with $\dprimeddprime(f\circ F)$; it is of course also
the ordinary chain rule applied twice.
\end{proof}

\begin{proposition}
\label{prop:ch09-pullback-functorial}
Pullback is functorial and commutes with $\dprime$, $\ddprime$, $J$, and wedge products.
Consequently the bigraded algebra and its differentials are independent of
affine coordinates.
\end{proposition}
\begin{proof}
The chain rule gives the assertions on generators, and the derivative of
$F$ is constant. Since functions together with the $\dprime y_i,\ddprime y_i$ generate
the algebra, the assertions extend to all forms. Applying this to inverse
affine coordinate changes proves coordinate independence.

We spell this out, since the argument is the model for every later
functoriality statement.  Let $F:U\to U'$ and $G:U'\to U''$ be restrictions
of affine maps, with linear parts $B$ and $C$.

\emph{Functoriality.}  On functions, $(G\circ F)^*a=a\circ G\circ F=F^*G^*a$.
On generators, $(G\circ F)^*\dprime z_k=\dprime(z_k\circ G\circ F)$, and this
equals $F^*\dprime(z_k\circ G)$ because $z_k\circ G$ is affine and, by the
computation in \Cref{lem:ch09-pullback-eleven}, $F^*$ applied to the
differential of an affine function is the differential of its composite with
$F$.  Since both $(G\circ F)^*$ and $F^*\circ G^*$ are multiplicative and
agree on generators of the algebra, they agree.  The identity map pulls back
to the identity, so $F\mapsto F^*$ is a contravariant functor.

\emph{Wedge products and $J$.}  Multiplicativity holds by construction.  For
the involution it suffices to check on generators:
$F^*J\dprime y_k=F^*\ddprime y_k=\ddprime(y_k\circ F)$ and
$JF^*\dprime y_k=J\dprime(y_k\circ F)=\ddprime(y_k\circ F)$, since
$J$ and $F^*$ are both algebra maps and $y_k\circ F$ is affine; the case of
$\ddprime y_k$ is symmetric, and on functions both sides are $a\circ F$.

\emph{Differentials.}  Let $\alpha=a\,\dprime y_K\wedge\ddprime y_L$.  Then
\[
 F^*\alpha=(a\circ F)\,F^*(\dprime y_K)\wedge F^*(\ddprime y_L),
\]
and the second and third factors have \emph{constant} coefficients, because
the linear part $B$ of $F$ is constant.  Applying
\eqref{eq:ch09-dprime-monomial},
\[
 \dprime F^*\alpha
 =\sum_i\frac{\partial (a\circ F)}{\partial x_i}\,
   \dprime x_i\wedge F^*(\dprime y_K)\wedge F^*(\ddprime y_L).
\]
The chain rule gives
$\partial(a\circ F)/\partial x_i=\sum_kB_{ki}\,(\partial a/\partial y_k)\circ F$,
so
\[
 \sum_i\frac{\partial (a\circ F)}{\partial x_i}\dprime x_i
 =\sum_k\Bigl(\frac{\partial a}{\partial y_k}\circ F\Bigr)
   \sum_iB_{ki}\dprime x_i
 =F^*\Bigl(\sum_k\frac{\partial a}{\partial y_k}\dprime y_k\Bigr),
\]
whence $\dprime F^*\alpha=F^*\dprime\alpha$.  The same computation with the
sign $(-1)^{|K|}$, or alternatively conjugation by $J$ using
\Cref{lem:ch09-J-conjugation} and the compatibility of $F^*$ with $J$ just
proved, gives $\ddprime F^*=F^*\ddprime$.

\emph{Coordinate independence.}  Let $x$ and $\tilde x$ be two affine
coordinate systems on $V$, and let $\Phi$ be the affine automorphism carrying
one to the other.  The constructions of \Cref{def:ch09-superform-space} and
\Cref{def:ch09-operators} performed in the two systems are intertwined by
$\Phi^*$, which is an isomorphism of bigraded algebras commuting with
$\dprime$, $\ddprime$ and $J$ by what precedes.  Hence the bigraded algebra,
its two differentials, and its involution depend only on $V$ and $U$, not on
the coordinates used to write them down.
\end{proof}

\begin{proposition}
\label{prop:ch09-pullback-top}
Let $F:V\to V$ be an affine self-map of an $n$-dimensional real vector space
with linear part $B$, and let $\tau_n$ be the top form attached to a basis as
in \Cref{lem:ch09-interlaced}.  Then
\[
 F^*\tau_n=\bigl(\det B\bigr)^2\,\tau_n .
\]
More generally, if $\alpha\in A^{n,n}(U)$ is written $\alpha=g\,\tau_n$, then
$F^*\alpha=(\det B)^2(g\circ F)\,\tau_n$.
\end{proposition}

\begin{proof}
By \Cref{lem:ch09-interlaced}, $\tau_n=c_n\,\lambda\wedge J\lambda$ with
$\lambda=\dprime x_1\wedge\cdots\wedge\dprime x_n$.  Now $F^*$ acts on the
primed block through $\bigwedge^n B^{\mathsf T}$, which on the
one-dimensional space $\bigwedge^nV^*$ is multiplication by $\det B$; hence
$F^*\lambda=(\det B)\lambda$.  Since $F^*$ commutes with $J$ by
\Cref{prop:ch09-pullback-functorial}, also $F^*(J\lambda)=(\det B)J\lambda$.
Multiplicativity gives
$F^*\tau_n=c_n(\det B)^2\lambda\wedge J\lambda=(\det B)^2\tau_n$.  The last
statement follows since $F^*(g\tau_n)=(g\circ F)F^*\tau_n$.
\end{proof}

The square in \Cref{prop:ch09-pullback-top} is the single most consequential
sign-free fact of the chapter.  It says that top-degree superforms transform
under affine maps by the \emph{square} of the Jacobian determinant, whereas
ordinary top-degree differential forms transform by the determinant itself.
Two consequences follow immediately.  First, positivity of a top form is
preserved by \emph{every} affine isomorphism, orientation-reversing ones
included: there is no orientation ambiguity in the positivity of a
Lagerberg top form.  Second, the naive integral
$\int_Ug\,\rd x$ attached to $\alpha=g\tau_n$ is \emph{not} invariant, since
the change of variables formula contributes only one factor $|\det B|$ while
the form contributes $(\det B)^2$.  The mismatch is a genuine feature, not a
defect: it is exactly what the calibrations of
\Cref{chap:tropical-integration} are designed to repair, and the surviving
factor $|\det B|$ is what becomes a lattice index in the tropical setting.

\begin{example}[A linear pullback computed in full]
\label{ex:ch09-pullback-explicit}
Let $F:\mathbf R^2\to\mathbf R^2$ be the linear map with matrix
\[
 B=\begin{pmatrix}1&1\\0&2\end{pmatrix},
 \qquad
 F(x_1,x_2)=(x_1+x_2,\;2x_2),
\]
and use coordinates $y_1,y_2$ on the target.  On generators,
\[
 \begin{aligned}
  F^*\dprime y_1&=\dprime x_1+\dprime x_2,
  &F^*\dprime y_2&=2\,\dprime x_2,\\
  F^*\ddprime y_1&=\ddprime x_1+\ddprime x_2,
  &F^*\ddprime y_2&=2\,\ddprime x_2 .
 \end{aligned}
\]
Take
\[
 \alpha=\dprime y_1\wedge\ddprime y_1
       +3\,\dprime y_2\wedge\ddprime y_2,
 \qquad
 A=\begin{pmatrix}1&0\\0&3\end{pmatrix}.
\]
Direct substitution gives
\[
 F^*\alpha
 =(\dprime x_1+\dprime x_2)\wedge(\ddprime x_1+\ddprime x_2)
  +12\,\dprime x_2\wedge\ddprime x_2,
\]
that is,
\[
 F^*\alpha
 =\dprime x_1\wedge\ddprime x_1
  +\dprime x_1\wedge\ddprime x_2
  +\dprime x_2\wedge\ddprime x_1
  +13\,\dprime x_2\wedge\ddprime x_2 .
\]
The coefficient matrix is $\left(\begin{smallmatrix}1&1\\1&13\end{smallmatrix}\right)$, and indeed
\[
 B^{\mathsf T}AB
 =\begin{pmatrix}1&0\\1&2\end{pmatrix}
  \begin{pmatrix}1&0\\0&3\end{pmatrix}
  \begin{pmatrix}1&1\\0&2\end{pmatrix}
 =\begin{pmatrix}1&1\\1&13\end{pmatrix},
\]
confirming \Cref{lem:ch09-pullback-eleven}.  Two checks are worth making.
The determinant is $\det(B^{\mathsf T}AB)=(\det B)^2\det A=4\cdot3=12$, so by
\Cref{ex:ch09-plane}
\[
 (F^*\alpha)^2=2\cdot12\,\tau_2=24\,\tau_2,
 \qquad\text{whereas}\qquad
 \alpha^2=2\det A\,\tau_2=6\,\tau_2 .
\]
Since $F^*$ is multiplicative, $(F^*\alpha)^2=F^*(\alpha^2)=6\,F^*\tau_2$,
and the two computations agree exactly because
$F^*\tau_2=(\det B)^2\tau_2=4\tau_2$, in agreement with
\Cref{prop:ch09-pullback-top}.  And $\alpha$ is positive, its matrix
being diagonal with nonnegative entries; so is $F^*\alpha$, whose matrix has
positive trace and positive determinant.  Note that $F$ is not an isometry
and does not preserve the Euclidean structure; positivity, however, is an
affine notion and survives.
\end{example}

\begin{example}[Restriction to an affine line]
\label{ex:ch09-restriction-line}
Let $F(t)=a+tv$ and let $f\in C^2(U')$.  Pulling back
$\dprimeddprime f$ gives
\[
 F^*(\dprimeddprime f)
 =v^{\mathsf T}(\operatorname{Hess}f)(a+tv)v\,
   \dprime t\wedge\ddprime t.
\]
Thus positivity of $\dprimeddprime f$ can be tested on affine lines.  This
calculation will reappear when psh functions are tested on analytic curves.
\end{example}

\begin{warning}{Why the map must be affine}
There is an evident algebraic pullback by an arbitrary smooth map if one
pulls back both labelled one-form factors by its derivative, but it does not
commute with both differentials.  Indeed, for a nonlinear function
$F:\mathbf R\to\mathbf R$, one has
$F^*(\dprimeddprime x)=0$ whereas $\dprimeddprime(F^*x)=F''\dprime x\wedge \ddprime x$.
The functorial calculus used in tropical geometry is therefore attached to
affine maps (and, globally, to the integer-affine transition maps of
tropical charts).
\end{warning}

The warning concerns functoriality of the differential calculus, not merely
notation.  A nonlinear change of variables produces second-derivative terms
even when the original form is zero.  Tropical transition maps are integral
affine precisely so that these terms never occur.

\section{Symmetry and the coefficient dictionary}
\label{sec:ch09-symmetry}

Positivity will be defined in \Cref{sec:ch09-positivity} by three conditions
which look, at first sight, like three unrelated pieces of formalism.  They
are not: each of them is a familiar condition on a symmetric matrix, once the
right matrix has been attached to a symmetric $(p,p)$-form.  This section
constructs that matrix and proves the one formula---the complementary
pairing---which converts wedge products of complementary bidegree into matrix
pairings.  Everything in the next section is then a computation with
symmetric matrices, and the classical counterexamples become explicit.

Throughout, $S_p$ denotes the set of $p$-element subsets of
$\{1,\ldots,n\}$, written as strictly increasing multi-indices; thus
$\#S_p=\binom np$ and $\dim A^{p,q}(V)=\#S_p\cdot\#S_q$.  For $I\in S_p$ we
write $\complement I\in S_{n-p}$ for the complementary multi-index, and
\[
 \varepsilon_I\in\{\pm1\}
\]
for the signature of the permutation of $\{1,\ldots,n\}$ obtained by listing
$I$ in increasing order and then $\complement I$ in increasing order.
Equivalently, $\varepsilon_I$ is defined by
\[
 \dprime x_I\wedge\dprime x_{\complement I}=\varepsilon_I\,\lambda,
 \qquad
 \lambda=\dprime x_1\wedge\cdots\wedge\dprime x_n .
\]

Put $c_p=(-1)^{p(p-1)/2}$. A $(p,p)$-form $\alpha$ is
\emph{symmetric} if $J\alpha=(-1)^p\alpha$.

\begin{proposition}
\label{prop:ch09-symmetry-matrix}
Let $\alpha=\sum_{I,J\in S_p}\alpha_{IJ}\,\dprime x_I\wedge\ddprime x_J$ be a
$(p,p)$-form.  Then:
\begin{enumerate}
\item $\alpha$ is symmetric if and only if $\alpha_{IJ}=\alpha_{JI}$ for all
  $I,J\in S_p$;
\item the condition in \textup{(1)} does not depend on the affine coordinates
  used to write $\alpha$;
\item symmetric forms are stable under products: if $\alpha$ and $\beta$ are
  symmetric of bidegrees $(p,p)$ and $(q,q)$, then $\alpha\wedge\beta$ is
  symmetric of bidegree $(p+q,p+q)$;
\item every $(n,n)$-form is symmetric, and so is $\dprimeddprime u$ for every
  $u\in C^2(U)$.
\end{enumerate}
\end{proposition}

\begin{proof}
(1) By \eqref{eq:ch09-J-coordinates} with $p=q$,
\[
 J\alpha=(-1)^{p^2}\sum_{I,J}\alpha_{IJ}\,\dprime x_J\wedge\ddprime x_I
        =(-1)^{p}\sum_{I,J}\alpha_{JI}\,\dprime x_I\wedge\ddprime x_J,
\]
where the second equality exchanges the names of the summation indices and
uses $(-1)^{p^2}=(-1)^p$.  Since the monomials
$\dprime x_I\wedge\ddprime x_J$ form a basis, $J\alpha=(-1)^p\alpha$ is
equivalent to $\alpha_{JI}=\alpha_{IJ}$ for all $I,J$.

(2) The condition $J\alpha=(-1)^p\alpha$ makes no reference to coordinates,
and $J$ is coordinate-independent by
\Cref{prop:ch09-pullback-functorial}.

(3) $J$ is an algebra automorphism, so
$J(\alpha\wedge\beta)=J\alpha\wedge J\beta=(-1)^{p+q}\alpha\wedge\beta$.

(4) The space $A^{n,n}$ is spanned by
$\dprime x_{\{1,\ldots,n\}}\wedge\ddprime x_{\{1,\ldots,n\}}$, whose single
coefficient array is trivially symmetric.  For $\dprimeddprime u$ the
coefficient matrix is the Hessian, symmetric by the equality of mixed partial
derivatives.
\end{proof}

Assertion (4) is the reason the whole positivity theory is set up for
symmetric forms: the objects it is meant to test---Hessians of potentials and
their products---are symmetric, and no information is lost.  Antisymmetric
forms are not merely uninteresting here, they are invisible: by
\Cref{lem:ch09-complement-pairing} below, a $(p,p)$-form with
$\alpha_{IJ}=-\alpha_{JI}$ pairs to zero with every symmetric
$(n-p,n-p)$-form.

\begin{definition}
\label{def:ch09-coefficient-matrix}
Fix a basis $x_1,\ldots,x_n$ of $V^*$ and let $\alpha\in A^{p,p}(V)$.  The
\emph{normalized coefficient matrix} of $\alpha$ is the array
$M(\alpha)=(\alpha_{IJ})_{I,J\in S_p}$ determined by
\[
 \alpha=c_p\sum_{I,J\in S_p}\alpha_{IJ}\,\dprime x_I\wedge\ddprime x_J,
 \qquad c_p=(-1)^{p(p-1)/2}.
\]
It is a symmetric matrix precisely when $\alpha$ is a symmetric form, and
$\alpha\mapsto M(\alpha)$ is a linear isomorphism from $A^{p,p}(V)$ onto the
space of $\binom np\times\binom np$ matrices.
\end{definition}

The factor $c_p$ in \Cref{def:ch09-coefficient-matrix} is a normalization and
nothing more; it is inserted so that the pairing formula below carries no
sign, and so that the generators of the positivity cones become rank-one
matrices with no sign attached.  For $p=1$ one has $c_1=1$ and the definition
reduces to the familiar coefficient matrix of a $(1,1)$-form; for $p=n$ one
has $M(\alpha)=g$ where $\alpha=g\,\tau_n$, since $c_n\lambda\wedge J\lambda
=\tau_n$ by \Cref{lem:ch09-interlaced}.

\begin{lemma}
\label{lem:ch09-complement-pairing}
Let $0\leq p\leq n$, $\alpha\in A^{p,p}(V)$ and
$\beta\in A^{n-p,n-p}(V)$.  Then
\[
 \alpha\wedge\beta
 =\Bigl(\sum_{I,J\in S_p}
    \varepsilon_I\varepsilon_J\,
    \alpha_{IJ}\,\beta_{\complement I,\complement J}\Bigr)\,\tau_n .
\]
Equivalently, defining the \emph{complementation operator}
\[
 (\star M)_{IJ}=\varepsilon_I\varepsilon_J\,M_{\complement I,\complement J},
 \qquad I,J\in S_p,
\]
which carries $S_{n-p}$-indexed matrices to $S_p$-indexed ones, one has
\[
 \alpha\wedge\beta=\bigl\langle M(\alpha),\star M(\beta)\bigr\rangle\,\tau_n,
 \qquad
 \langle A,B\rangle:=\sum_{I,J}A_{IJ}B_{IJ}.
\]
The operator $\star$ preserves symmetry of matrices and satisfies
$\star\star=\operatorname{id}$.
\end{lemma}

\begin{proof}
Write $\alpha=c_p\sum_{I,J}\alpha_{IJ}\dprime x_I\wedge\ddprime x_J$ and
$\beta=c_{n-p}\sum_{K,L}\beta_{KL}\dprime x_K\wedge\ddprime x_L$ with
$I,J\in S_p$ and $K,L\in S_{n-p}$.  The basic product is
\[
 \bigl(\dprime x_I\wedge\ddprime x_J\bigr)\wedge
 \bigl(\dprime x_K\wedge\ddprime x_L\bigr)
 =(-1)^{p(n-p)}\,
   \dprime x_I\wedge\dprime x_K\wedge\ddprime x_J\wedge\ddprime x_L,
\]
by \Cref{prop:ch09-wedge-signs}\,(1) applied with $q=p$ and $p'=n-p$: the
double-primed block $\ddprime x_J$, of degree $p$, is moved past the primed
block $\dprime x_K$, of degree $n-p$.  Now
$\dprime x_I\wedge\dprime x_K$ vanishes unless $K=\complement I$, in which
case it equals $\varepsilon_I\lambda$; similarly
$\ddprime x_J\wedge\ddprime x_L=\varepsilon_J\,J\lambda$ if
$L=\complement J$, and vanishes otherwise.  Hence only the terms with
$K=\complement I$ and $L=\complement J$ survive, and
\[
 \alpha\wedge\beta
 =c_pc_{n-p}(-1)^{p(n-p)}
   \sum_{I,J}\varepsilon_I\varepsilon_J\,
   \alpha_{IJ}\beta_{\complement I,\complement J}\;
   \lambda\wedge J\lambda .
\]
By \Cref{lem:ch09-interlaced} we have $\lambda\wedge J\lambda=c_n\tau_n$ and
$c_pc_{n-p}(-1)^{p(n-p)}=c_n$, so the total constant is $c_n^2=1$.  This is
the asserted formula.

That $\star$ preserves symmetry is immediate from its definition.  For the
last claim, note that $\varepsilon_I\varepsilon_{\complement I}
=(-1)^{p(n-p)}$: indeed
$\dprime x_I\wedge\dprime x_{\complement I}=\varepsilon_I\lambda$ and
$\dprime x_{\complement I}\wedge\dprime x_I=\varepsilon_{\complement I}\lambda$,
while interchanging the two blocks costs $(-1)^{p(n-p)}$.  Therefore
\[
 (\star\star M)_{IJ}
 =\varepsilon_I\varepsilon_J(\star M)_{\complement I,\complement J}
 =\varepsilon_I\varepsilon_J\varepsilon_{\complement I}
   \varepsilon_{\complement J}M_{IJ}
 =(-1)^{2p(n-p)}M_{IJ}=M_{IJ}.
\]
\end{proof}

\begin{warning}{The pairing is sign-free only with the normalization $c_p$}
If the factor $c_p$ is dropped from
\Cref{def:ch09-coefficient-matrix}, the pairing formula acquires the sign
$c_pc_{n-p}(-1)^{p(n-p)}c_n$, which equals $1$ only by the identity of
\Cref{lem:ch09-interlaced}.  This is the sole reason for the constants
$c_p$ recorded as \eqref{eq:appJ-cp} in \Cref{sec:appJ-cld}.  In particular a
strongly positive generator carries a single factor $c_p$ in total, not one
factor per rank-one term; inserting one per term would double the sign in
even bidegree and reverse the inclusions between the positivity cones.
\end{warning}

We record the vectorial form of $\star$, which is what the positivity
criteria use.  For $b=(b_K)_{K\in S_{n-p}}$ set
\[
 (\star b)_I=\varepsilon_I\,b_{\complement I},\qquad I\in S_p .
\]
Then $\star(bb^{\mathsf T})=(\star b)(\star b)^{\mathsf T}$, directly from the
definitions.  We identify a vector $(a_I)_{I\in S_p}$ with the element
$\sum_Ia_I\,\dprime x_I$ of $A^{p,0}(V)=\bigwedge^pV^*$ and call it
\emph{decomposable} if that element is a product
$\eta_1\wedge\cdots\wedge\eta_p$ of elements of $V^*$; the zero vector is
decomposable.  Write $\Gamma_p\subset\bigwedge^pV^*$ for the set of
decomposable elements.  It is a closed cone, being the zero locus of the
finitely many quadratic Pl\"ucker relations, or, more elementarily, the image
of the continuous proper-on-spheres map
$(V^*)^p\to\bigwedge^pV^*$ restricted to bounded sets; we shall only use that
it is closed and stable under scaling, both of which are clear from the
description as an image of a map, together with the fact that a limit of
decomposable elements is decomposable.

\begin{lemma}
\label{lem:ch09-star-decomposable}
For $a\in\bigwedge^pV^*$ and $b\in\bigwedge^{n-p}V^*$ one has
\[
 a\wedge b=\langle a,\star b\rangle\,\lambda,
\]
where $\langle\cdot,\cdot\rangle$ is the inner product for which the basis
$(\dprime x_I)_{I\in S_p}$ is orthonormal.  Moreover $\star$ restricts to a
bijection
\[
 \star:\Gamma_{n-p}\xrightarrow{\ \sim\ }\Gamma_p .
\]
\end{lemma}

\begin{proof}
For the first formula, expand
$a\wedge b=\sum_{I,K}a_Ib_K\,\dprime x_I\wedge\dprime x_K$; only $K=\complement I$
contributes, giving $\sum_Ia_I\varepsilon_Ib_{\complement I}\lambda
=\langle a,\star b\rangle\lambda$.

Let now $b\in\Gamma_{n-p}$ be nonzero, say
$b=\beta_{p+1}\wedge\cdots\wedge\beta_n$ with $\beta_i\in V^*$ linearly
independent.  Apply the Gram--Schmidt process to
$\beta_{p+1},\ldots,\beta_n$ and complete to an orthonormal basis
$\gamma_1,\ldots,\gamma_n$ of $V^*$ such that
$\gamma_{p+1},\ldots,\gamma_n$ span the same subspace as
$\beta_{p+1},\ldots,\beta_n$.  Since the base change between the two families
of that subspace is triangular with nonzero determinant, there is $c\neq0$
with
\[
 b=c\,\gamma_{p+1}\wedge\cdots\wedge\gamma_n .
\]
Recall that if $(\gamma_i)$ is orthonormal in $V^*$, then the products
$\gamma_K$, $K\in S_p$, form an orthonormal basis of $\bigwedge^pV^*$: this is
the Gram determinant formula
$\langle\gamma_K,\gamma_L\rangle=\det(\langle\gamma_{k_i},\gamma_{l_j}\rangle)
=\delta_{KL}$.  Write $d=\pm1$ for the determinant of the base change from
$(\dprime x_i)$ to $(\gamma_i)$, so that
$\gamma_1\wedge\cdots\wedge\gamma_n=d\,\lambda$.  Then for every
$a\in\bigwedge^pV^*$,
\[
 a\wedge b
 =c\,a\wedge\gamma_{p+1}\wedge\cdots\wedge\gamma_n
 =c\,\bigl\langle a,\gamma_1\wedge\cdots\wedge\gamma_p\bigr\rangle\,
   \gamma_1\wedge\cdots\wedge\gamma_n
 =cd\,\bigl\langle a,\gamma_1\wedge\cdots\wedge\gamma_p\bigr\rangle\lambda,
\]
the middle equality because, expanding $a$ in the orthonormal basis
$(\gamma_K)$, only the coefficient along $\gamma_1\wedge\cdots\wedge\gamma_p$
survives the wedge with $\gamma_{p+1}\wedge\cdots\wedge\gamma_n$.  Comparing
with the first formula and using that the inner product is nondegenerate, we
conclude
\[
 \star b=cd\,\gamma_1\wedge\cdots\wedge\gamma_p\in\Gamma_p .
\]
Thus $\star(\Gamma_{n-p})\subset\Gamma_p$.  The same argument with $p$ and
$n-p$ exchanged gives $\star(\Gamma_p)\subset\Gamma_{n-p}$ for the
complementation operator in the other direction, and the two are inverse to
each other up to the sign computed in
\Cref{lem:ch09-complement-pairing}, whence the asserted bijection.
\end{proof}

\section{Positivity}
\label{sec:ch09-positivity}

\begin{definition}
\label{def:ch09-positivity}
A symmetric $(p,p)$-form is \emph{strongly positive} if locally it is a finite
sum with nonnegative smooth coefficients of forms
\[
\eta_1\wedge J\eta_1\wedge\cdots\wedge
                 \eta_p\wedge J\eta_p,\qquad \eta_i\in A^{1,0}.
\]
It is \emph{positive} if it is locally a nonnegative sum of
$c_p\eta\wedge J\eta$ with $\eta\in A^{p,0}$. It is \emph{weakly positive} if
$\alpha\wedge\beta$ is a nonnegative top form for every strongly positive
$(n-p,n-p)$-form $\beta$.
\end{definition}

Here the Lagerberg orientation associated to an affine basis is
\[
\tau_n=\dprime x_1\wedge \ddprime x_1\wedge\cdots\wedge \dprime x_n\wedge \ddprime x_n.
\]
A top form is nonnegative when it is $f\tau_n$ with $f\geq0$.
By \Cref{prop:ch09-pullback-top} this notion is invariant under every affine
isomorphism, orientation-reversing ones included, so no orientation of $V$
has been chosen; $\tau_n$ depends on the basis only through the square of the
determinant of the base change.

\paragraph{Positivity is tested pointwise.}
Throughout this section we first define three convex cones inside the
finite-dimensional vector space $A^{p,p}(V)$ and then declare a superform on
$U$ to be strongly positive, positive, or weakly positive when its value at
each point of $U$ lies in the corresponding cone.  A superform which is
locally a finite sum with nonnegative smooth coefficients of the generators
listed in \Cref{def:ch09-positivity} clearly has this property, so the reading
adopted here is the weaker and safer one.  The converse question---whether a
pointwise decomposition into generators can be chosen to depend smoothly on
the point---is a delicate matter of real analysis; it is never used in this
book and is not asserted anywhere in it.  Every criterion proved below is
pointwise, and every application is pointwise.

\begin{definition}
\label{def:ch09-cones}
For $0\leq p\leq n$ define subsets of the space
$A^{p,p}_{\mathrm{sym}}(V)$ of symmetric $(p,p)$-forms on $V$ by
\[
\begin{aligned}
 \mathrm{SP}_p(V)
 &=\Bigl\{\textstyle\sum_{s}a_s\,
    \eta_{1,s}\wedge J\eta_{1,s}\wedge\cdots\wedge
    \eta_{p,s}\wedge J\eta_{p,s}
   \;:\;a_s\geq0,\ \eta_{i,s}\in A^{1,0}(V)\Bigr\},\\
 \mathrm{P}_p(V)
 &=\Bigl\{\textstyle\sum_{s}a_s\,c_p\,\eta_s\wedge J\eta_s
   \;:\;a_s\geq0,\ \eta_s\in A^{p,0}(V)\Bigr\},\\
 \mathrm{WP}_p(V)
 &=\bigl\{\alpha\in A^{p,p}_{\mathrm{sym}}(V)\;:\;
   \alpha\wedge\beta\in\mathbf R_{\geq0}\,\tau_n
   \text{ for all }\beta\in\mathrm{SP}_{n-p}(V)\bigr\},
\end{aligned}
\]
all sums being finite.  These are the \emph{strongly positive}, the
\emph{positive}, and the \emph{weakly positive} cone in bidegree $(p,p)$.
\end{definition}

That $\mathrm{SP}_p(V)$ and $\mathrm{P}_p(V)$ consist of symmetric forms is
part of the assertion of the next lemma; that all three are convex cones is
immediate from the definitions, the last because the defining conditions are
linear inequalities.

\begin{lemma}
\label{lem:ch09-generators}
Let $\eta_1,\ldots,\eta_p\in A^{1,0}(V)$ and put
$\eta=\eta_1\wedge\cdots\wedge\eta_p\in A^{p,0}(V)$.  Then
\[
 \eta_1\wedge J\eta_1\wedge\cdots\wedge\eta_p\wedge J\eta_p
 =c_p\,\eta\wedge J\eta .
\]
More generally, for $\eta\in A^{a,0}(V)$ and $\zeta\in A^{b,0}(V)$,
\[
 \bigl(c_a\,\eta\wedge J\eta\bigr)\wedge
 \bigl(c_b\,\zeta\wedge J\zeta\bigr)
 =c_{a+b}\,(\eta\wedge\zeta)\wedge J(\eta\wedge\zeta).
\]
Every form $c_a\eta\wedge J\eta$ with $\eta\in A^{a,0}(V)$ is symmetric, and
its normalized coefficient matrix is $vv^{\mathsf T}$, where
$v=(\eta_I)_{I\in S_a}$ is the coefficient vector of $\eta$.  Consequently
$\mathrm{SP}_p(V)\subset\mathrm{P}_p(V)$, and both consist of symmetric
forms.
\end{lemma}

\begin{proof}
For the second formula, move the block $J\eta$, of total degree $a$, past the
block $\zeta$, of total degree $b$, which costs $(-1)^{ab}$, and use that $J$
is an algebra automorphism:
\[
 \eta\wedge J\eta\wedge\zeta\wedge J\zeta
 =(-1)^{ab}\,\eta\wedge\zeta\wedge J\eta\wedge J\zeta
 =(-1)^{ab}\,(\eta\wedge\zeta)\wedge J(\eta\wedge\zeta).
\]
Multiplying by $c_ac_b$ and invoking the identity
$c_ac_b(-1)^{ab}=c_{a+b}$ of \Cref{lem:ch09-interlaced} gives the claim.  The
first formula follows by induction on $p$: it is trivial for $p=1$ since
$c_1=1$, and the inductive step is the second formula with $a=p-1$, $b=1$.

For the coefficient matrix, write $\eta=\sum_I\eta_I\dprime x_I$ with
$I\in S_a$.  Then $J\eta=\sum_J\eta_J\ddprime x_J$, so
\[
 c_a\,\eta\wedge J\eta
 =c_a\sum_{I,J}\eta_I\eta_J\,\dprime x_I\wedge\ddprime x_J,
\]
which by \Cref{def:ch09-coefficient-matrix} means
$M(c_a\eta\wedge J\eta)=(\eta_I\eta_J)_{I,J}=vv^{\mathsf T}$.  This matrix is
symmetric, so the form is symmetric by
\Cref{prop:ch09-symmetry-matrix}; sums of such forms with nonnegative
coefficients are then symmetric as well.  Finally, the first formula
identifies each generator of $\mathrm{SP}_p(V)$ with a generator of
$\mathrm{P}_p(V)$, whence the inclusion.
\end{proof}

The lemma already isolates the essential difference between the two cones:
$\mathrm{SP}_p$ is generated by the matrices $vv^{\mathsf T}$ with $v$
\emph{decomposable}, while $\mathrm{P}_p$ is generated by all of them.  The
next proposition turns this observation into a complete dictionary.

\begin{proposition}
\label{prop:ch09-positivity-dictionary}
Fix a basis of $V^*$, let $\alpha\in A^{p,p}_{\mathrm{sym}}(V)$ and let
$M=M(\alpha)$ be its normalized coefficient matrix, a symmetric matrix
indexed by $S_p$.  Recall $\Gamma_p\subset\bigwedge^pV^*$ from
\Cref{sec:ch09-symmetry}, identified with a subset of $\mathbf R^{S_p}$ by
means of the coefficient vectors.  Then:
\begin{enumerate}
\item $\alpha\in\mathrm{P}_p(V)$ if and only if $M$ is positive
  semidefinite;
\item $\alpha\in\mathrm{SP}_p(V)$ if and only if
  $M=\sum_sa_sv_sv_s^{\mathsf T}$ for finitely many $a_s\geq0$ and
  $v_s\in\Gamma_p$;
\item $\alpha\in\mathrm{WP}_p(V)$ if and only if $w^{\mathsf T}Mw\geq0$ for
  every $w\in\Gamma_p$.
\end{enumerate}
\end{proposition}

\begin{proof}
By \Cref{lem:ch09-generators} the map $\alpha\mapsto M(\alpha)$ carries the
generators $c_p\eta\wedge J\eta$ of $\mathrm{P}_p(V)$ to the matrices
$vv^{\mathsf T}$ with $v\in\mathbf R^{S_p}$ arbitrary, and the generators of
$\mathrm{SP}_p(V)$ to those with $v\in\Gamma_p$.  Since $M$ is a linear
isomorphism onto the space of matrices, it carries $\mathrm{P}_p(V)$ onto
the convex cone generated by all $vv^{\mathsf T}$ and $\mathrm{SP}_p(V)$ onto
the convex cone generated by the $vv^{\mathsf T}$ with $v\in\Gamma_p$.  This
is (2), and it gives (1) once one recalls that the convex cone generated by
all rank-one squares is exactly the cone of positive semidefinite matrices:
one inclusion is clear, and the other is the spectral theorem, which writes a
positive semidefinite $M$ as $\sum_k\mu_kw_kw_k^{\mathsf T}$ with
$\mu_k\geq0$.

For (3), let $\beta\in\mathrm{SP}_{n-p}(V)$ and write, by (2) applied in
bidegree $(n-p,n-p)$, $M(\beta)=\sum_sb_sb_s^{\mathsf T}$ with
$b_s\in\Gamma_{n-p}$; the nonnegative coefficients have been absorbed into
the vectors.  By \Cref{lem:ch09-complement-pairing},
\[
 \alpha\wedge\beta
 =\bigl\langle M,\star M(\beta)\bigr\rangle\tau_n
 =\sum_s\bigl\langle M,(\star b_s)(\star b_s)^{\mathsf T}\bigr\rangle\tau_n
 =\sum_s(\star b_s)^{\mathsf T}M(\star b_s)\;\tau_n,
\]
where we used $\star(bb^{\mathsf T})=(\star b)(\star b)^{\mathsf T}$ and the
identity $\langle M,ww^{\mathsf T}\rangle=w^{\mathsf T}Mw$.  By
\Cref{lem:ch09-star-decomposable}, $\star$ maps $\Gamma_{n-p}$ bijectively
onto $\Gamma_p$.  Hence $\alpha\wedge\beta\in\mathbf R_{\geq0}\tau_n$ for all
$\beta\in\mathrm{SP}_{n-p}(V)$ if and only if $w^{\mathsf T}Mw\geq0$ for all
$w\in\Gamma_p$, which is (3).
\end{proof}

\begin{proposition}
\label{prop:ch09-positivity-inclusions}
Strong positivity implies positivity, and positivity implies weak positivity.
All three cones are convex, closed under multiplication by nonnegative
functions and under pullback. They coincide for $p=0,1,n-1,n$.
\end{proposition}
\begin{proof}
A decomposable $p$-form
$\eta=\eta_1\wedge\cdots\wedge\eta_p$ satisfies
$c_p\eta\wedge J\eta=
\eta_1\wedge J\eta_1\wedge\cdots\wedge\eta_p\wedge J\eta_p$ after the
signs are collected, proving the first inclusion. If $\alpha=c_p\eta\wedge
J\eta$ and $\beta$ is a strongly positive complementary form, their wedge is
a square of a determinant times the positive oriented volume form, proving
the second. The stability statements follow from the definitions. The extreme
bidegrees reduce to nonnegative functions, while for $p=1$ positivity is
positive semidefiniteness of a symmetric matrix; duality gives $p=n-1$.

We now give the details.  The first inclusion is
\Cref{lem:ch09-generators}.  For the second, let $\alpha=c_p\eta\wedge J\eta$
with $\eta\in A^{p,0}(V)$ and let
$\beta=c_{n-p}\zeta\wedge J\zeta$ with $\zeta\in\Gamma_{n-p}$ be a generator
of $\mathrm{SP}_{n-p}(V)$, which is legitimate by
\Cref{lem:ch09-generators}.  By the second formula of that lemma,
\[
 \alpha\wedge\beta
 =c_n\,(\eta\wedge\zeta)\wedge J(\eta\wedge\zeta),
\]
and $\eta\wedge\zeta\in\bigwedge^nV^*$ is a multiple $t\lambda$ of the
generator, so
$\alpha\wedge\beta=c_nt^2\lambda\wedge J\lambda=t^2\tau_n$ by
\Cref{lem:ch09-interlaced}.  This is a nonnegative top form; passing to
nonnegative combinations gives $\mathrm{P}_p(V)\subset\mathrm{WP}_p(V)$.
Alternatively, and more quickly: a positive semidefinite matrix satisfies
$w^{\mathsf T}Mw\geq0$ for every $w$, in particular for decomposable $w$, so
the inclusion is immediate from
\Cref{prop:ch09-positivity-dictionary}.

Convexity and stability under multiplication by a nonnegative function are
clear from \Cref{def:ch09-cones}, since each cone is defined by nonnegative
combinations or by linear inequalities, and since multiplying a generator by
$a\geq0$ produces a generator.

For stability under pullback, let $F:V\to V'$ be affine with linear part $B$
and let $\alpha\in A^{p,p}(V')$.  Since $F^*$ is multiplicative and commutes
with $J$ by \Cref{prop:ch09-pullback-functorial},
\[
 F^*\bigl(\eta_1\wedge J\eta_1\wedge\cdots\wedge\eta_p\wedge J\eta_p\bigr)
 =(F^*\eta_1)\wedge J(F^*\eta_1)\wedge\cdots
   \wedge(F^*\eta_p)\wedge J(F^*\eta_p),
\]
so $F^*$ maps generators of $\mathrm{SP}_p(V')$ to generators of
$\mathrm{SP}_p(V)$, and similarly for $\mathrm{P}_p$; both cones are
therefore preserved.  For $\mathrm{WP}_p$ we use the coefficient matrices.
Writing $C$ for the matrix of $\bigwedge^pB^{\mathsf T}$, whose entries are
the $p\times p$ minors $C_{I,K}=\det\bigl(B_{K,I}\bigr)$, the identity
$F^*\dprime y_K=\sum_IC_{I,K}\dprime x_I$ gives
\[
 M(F^*\alpha)=C\,M(\alpha)\,C^{\mathsf T},
\]
exactly as in \Cref{lem:ch09-pullback-eleven}, of which this is the case
$p=1$.  Now $C^{\mathsf T}$ maps $\Gamma_p(V)$ into $\Gamma_p(V')$, being the
matrix of $\bigwedge^pB$ in the corresponding bases and hence sending
$v_1\wedge\cdots\wedge v_p$ to $Bv_1\wedge\cdots\wedge Bv_p$.  Therefore, for
$w\in\Gamma_p(V)$,
\[
 w^{\mathsf T}M(F^*\alpha)w
 =\bigl(C^{\mathsf T}w\bigr)^{\mathsf T}M(\alpha)\bigl(C^{\mathsf T}w\bigr)
 \geq0
\]
whenever $\alpha$ is weakly positive, by
\Cref{prop:ch09-positivity-dictionary}\,(3).  Hence $F^*\alpha$ is weakly
positive.

Finally, the coincidence in the four extreme bidegrees.  By
\Cref{prop:ch09-positivity-dictionary} it suffices to prove that
$\Gamma_p=\bigwedge^pV^*$ for $p\in\{0,1,n-1,n\}$: if every vector is
decomposable, then the three conditions of that proposition all reduce to
positive semidefiniteness of $M$.  For $p=0$ and $p=1$ this is trivial, the
spaces being $\mathbf R$ and $V^*$.  For $p=n$ the space
$\bigwedge^nV^*$ is one-dimensional, spanned by the decomposable element
$\lambda$, and multiples of decomposable elements are decomposable.  For
$p=n-1$, apply \Cref{lem:ch09-star-decomposable}: the operator $\star$ is a
linear bijection of $\bigwedge^1V^*$ onto $\bigwedge^{n-1}V^*$ carrying
$\Gamma_1$ onto $\Gamma_{n-1}$, and $\Gamma_1=\bigwedge^1V^*$; hence
$\Gamma_{n-1}=\bigwedge^{n-1}V^*$.
\end{proof}

The proof identifies precisely what makes the three notions differ, and it is
worth stating in one sentence: \emph{the three cones coincide in bidegree
$(p,p)$ as soon as every $p$-covector is decomposable}.  Every $p$-covector
on an $n$-dimensional space is decomposable when $p\leq1$ or $p\geq n-1$, as
the proof just showed, so the first bidegree in which anything can go wrong
is $p=2$ on a space of dimension $n=4$.  It does go wrong there, as
\Cref{ex:ch09-weak-not-positive,ex:ch09-positive-not-strong} show by explicit
computation.  In particular all three notions agree in every bidegree when
$n\leq3$, which is why the distinction is invisible on curves and surfaces
and why it is so often mislaid.

\begin{proposition}
\label{prop:ch09-positivity-products}
Let $\alpha\in A^{p,p}(V)$ and $\beta\in A^{q,q}(V)$ with $p+q\leq n$.
\begin{enumerate}
\item If $\alpha$ and $\beta$ are strongly positive, so is
  $\alpha\wedge\beta$.
\item If $\alpha$ and $\beta$ are positive, so is $\alpha\wedge\beta$.
\item If $\alpha$ is weakly positive and $\beta$ is strongly positive, then
  $\alpha\wedge\beta$ is weakly positive.
\end{enumerate}
In particular the product of finitely many strongly positive forms is
strongly positive, and a product of weakly positive forms of which at most
one fails to be strongly positive is weakly positive.
\end{proposition}

\begin{proof}
(1)  By bilinearity we may assume $\alpha$ and $\beta$ are generators, say
$\alpha=\eta_1\wedge J\eta_1\wedge\cdots\wedge\eta_p\wedge J\eta_p$ and
$\beta=\zeta_1\wedge J\zeta_1\wedge\cdots\wedge\zeta_q\wedge J\zeta_q$.  Each
factor $\eta_i\wedge J\eta_i$ has total degree two, hence commutes with
everything by \Cref{prop:ch09-wedge-signs}\,(2), and the concatenated product
is again a generator, now with $p+q$ factors.

(2)  Again reduce to generators $\alpha=c_p\eta\wedge J\eta$ and
$\beta=c_q\zeta\wedge J\zeta$; the second formula of
\Cref{lem:ch09-generators} gives
$\alpha\wedge\beta=c_{p+q}(\eta\wedge\zeta)\wedge J(\eta\wedge\zeta)$, a
generator of $\mathrm{P}_{p+q}(V)$.

(3)  The form $\alpha\wedge\beta$ is symmetric by
\Cref{prop:ch09-symmetry-matrix}\,(3).  Let
$\gamma\in\mathrm{SP}_{n-p-q}(V)$.  By (1), $\beta\wedge\gamma$ lies in
$\mathrm{SP}_{n-p}(V)$, so
\[
 (\alpha\wedge\beta)\wedge\gamma=\alpha\wedge(\beta\wedge\gamma)
 \in\mathbf R_{\geq0}\tau_n
\]
by weak positivity of $\alpha$.  Hence $\alpha\wedge\beta$ is weakly
positive.  The final sentence follows by induction, using (1) to combine the
strongly positive factors first.
\end{proof}

\begin{warning}{Products of weakly positive forms}
Assertion (3) of \Cref{prop:ch09-positivity-products} is sharp: the product of
two weakly positive forms need not be weakly positive, and the product of two
positive forms is positive but need not be strongly positive.  The reason is
visible in the proof: only the strongly positive factor can be moved across
the pairing.  In the applications this is the reason that mixed
Monge--Amp\`ere products are formed with at most one factor of doubtful
positivity, and it is one of the sources of the hypotheses in
\Cref{chap:cld-bedford-taylor}.
\end{warning}

\begin{theorem}
\label{thm:ch09-positivity-duality}
Under the pairing
\[
 A^{p,p}_{\mathrm{sym}}(V)\times A^{n-p,n-p}_{\mathrm{sym}}(V)
 \longrightarrow\mathbf R,
 \qquad
 (\alpha,\beta)\longmapsto
 \langle\alpha,\beta\rangle\ \text{ defined by }\
 \alpha\wedge\beta=\langle\alpha,\beta\rangle\,\tau_n,
\]
which is a perfect pairing of finite-dimensional vector spaces, the three
cones satisfy
\[
 \mathrm{WP}_p(V)=\mathrm{SP}_{n-p}(V)^\vee,
 \qquad
 \mathrm{SP}_p(V)=\mathrm{WP}_{n-p}(V)^\vee,
 \qquad
 \mathrm{P}_p(V)=\mathrm{P}_{n-p}(V)^\vee,
\]
where $C^\vee$ denotes the cone of elements pairing nonnegatively with every
element of $C$.  In particular the positive cone is self-dual, and the two
extreme cones are dual to each other.
\end{theorem}

\begin{proof}
\emph{The pairing is perfect.}  By
\Cref{lem:ch09-complement-pairing},
$\langle\alpha,\beta\rangle=\langle M(\alpha),\star M(\beta)\rangle$ with
$\langle A,B\rangle=\sum_{I,J}A_{IJ}B_{IJ}$.  The latter is a positive
definite inner product on the space of matrices, $\star$ is a linear
involution, and $M$ is a linear isomorphism; so the pairing is nondegenerate,
and both spaces have the same dimension $\binom np\bigl(\binom np+1\bigr)/2$.

\emph{First identity.}  This is \Cref{def:ch09-cones} restated: a symmetric
form $\alpha$ is weakly positive precisely when it pairs nonnegatively with
every strongly positive form of complementary bidegree.

\emph{The cone $\mathrm{SP}_p(V)$ is closed.}  Let
$K=\{vv^{\mathsf T}:v\in\Gamma_p,\ |v|=1\}$, a subset of the space of
symmetric matrices.  It is compact, being the image, under a continuous map,
of the intersection of $\Gamma_p$ with the unit sphere of
$\mathbf R^{S_p}$; that intersection is compact because $\Gamma_p$ is closed
and the sphere is compact.  Every element of $K$ has trace
$|v|^2=1$, so the convex hull $\operatorname{conv}(K)$, which is compact by
Carath\'eodory's theorem in finite dimension, lies in the affine hyperplane
$\{\operatorname{tr}=1\}$ and in particular does not contain $0$.  Now
$M(\mathrm{SP}_p(V))$ is the set of nonnegative multiples of elements of
$\operatorname{conv}(K)$.  If $t_j\kappa_j\to A$ with $t_j\geq0$ and
$\kappa_j\in\operatorname{conv}(K)$, then, after passing to a subnet,
$\kappa_j\to\kappa\in\operatorname{conv}(K)$; since
$\operatorname{tr}(t_j\kappa_j)=t_j$ converges, $t_j\to t\geq0$ and
$A=t\kappa$.  Hence $M(\mathrm{SP}_p(V))$ is closed, and so is
$\mathrm{SP}_p(V)$.

\emph{Second identity.}  For a closed convex cone $C$ in a finite-dimensional
vector space paired perfectly with another, one has $C^{\vee\vee}=C$: the
inclusion $C\subset C^{\vee\vee}$ is formal, and if $x\notin C$ then, by the
separation theorem for a point and a closed convex set, there is a functional
$\xi$ with $\langle\xi,x\rangle<\inf_{c\in C}\langle\xi,c\rangle$; the
infimum over a cone is $0$ or $-\infty$, and it is finite, hence $0$, so
$\xi\in C^\vee$ and $\langle\xi,x\rangle<0$, whence $x\notin C^{\vee\vee}$.
See \cite[Theorem~14.1]{Roc70} for this statement in the general form.
Applying it to the closed convex cone $\mathrm{SP}_p(V)$ and using the first
identity gives
\[
 \mathrm{WP}_{n-p}(V)^\vee
 =\bigl(\mathrm{SP}_{p}(V)^\vee\bigr)^\vee
 =\mathrm{SP}_p(V).
\]

\emph{Third identity.}  By
\Cref{prop:ch09-positivity-dictionary}\,(1), $\mathrm{P}_p(V)$ corresponds
under $M$ to the cone of positive semidefinite matrices, which is self-dual
for $\langle\cdot,\cdot\rangle$: if $A=\sum_sv_sv_s^{\mathsf T}$ and $B$ are
positive semidefinite then
$\langle A,B\rangle=\sum_sv_s^{\mathsf T}Bv_s\geq0$, and conversely, taking
$B=ww^{\mathsf T}$ in $\langle A,B\rangle\geq0$ gives
$w^{\mathsf T}Aw\geq0$ for all $w$.  Since
$\langle\alpha,\beta\rangle=\langle M(\alpha),\star M(\beta)\rangle$ and
$\star$ maps positive semidefinite matrices to positive semidefinite matrices
---because $\star(bb^{\mathsf T})=(\star b)(\star b)^{\mathsf T}$ and $\star$
is a linear bijection on vectors---the assertion follows.
\end{proof}

The duality theorem is the reason the middle cone is the good one for many
purposes: it is the unique self-dual member of the triple, and, as
\Cref{prop:ch09-positivity-products}\,(2) shows, it is stable under products.
The extreme cones trade places under duality, which is also why statements
about weakly positive forms are so often proved by testing against strongly
positive ones rather than directly.

\begin{proposition}
\label{prop:superform-matrix-wedge}
Let
\[
 \alpha=\sum_{i,j=1}^n a_{ij}\,
                    \dprime x_i\wedge\ddprime x_j
\]
be a symmetric $(1,1)$-form. Then $A=(a_{ij})$ is symmetric,
$\alpha$ is positive if and only if $A$ is positive semidefinite, and
\[
 \alpha^n=n!\det(A)\,\tau_n.
\]
More generally, the coefficient of
$\alpha_1\wedge\cdots\wedge\alpha_n$ relative to $\tau_n$ is the mixed
discriminant of their coefficient matrices, with the normalization for which
the diagonal value is $n!\det(A)$.
\end{proposition}
\begin{proof}
Applying $J$ to a basis term and moving the double-primed factor past the
primed factor shows that the symmetry condition is $a_{ij}=a_{ji}$.  Every
positive semidefinite matrix is a finite sum of rank-one matrices
$vv^{\mathsf T}$.  The corresponding form is a sum of
$\eta\wedge J\eta$ with $\eta=\sum_i v_i\dprime x_i$, proving the positivity
criterion.  Conversely, each such rank-one summand is positive
semidefinite.  Expanding $\alpha^n$, only terms using every primed and every
double-primed index survive.  The alternating sum is $n!\det(A)$; the sign
needed to group the two labelled blocks is exactly cancelled by the sign
relating the grouped and interlaced orientations.  Polarization gives the
mixed statement.

We now supply the two computations in full.

\emph{Symmetry and positivity.}  These are the cases $p=1$ of
\Cref{prop:ch09-symmetry-matrix} and
\Cref{prop:ch09-positivity-dictionary}, since $c_1=1$ and every element of
$\bigwedge^1V^*$ is decomposable; in this bidegree all three positivity
notions agree by \Cref{prop:ch09-positivity-inclusions}.

\emph{The top power.}  Write $e_i=\dprime x_i$ and $f_j=\ddprime x_j$, so
that $\alpha=\sum_{i,j}a_{ij}e_i\wedge f_j$.  Each summand has total degree
two and therefore commutes with every other, so
\[
 \alpha^n
 =\sum_{i_1,\ldots,i_n}\ \sum_{j_1,\ldots,j_n}
   a_{i_1j_1}\cdots a_{i_nj_n}\,
   (e_{i_1}\wedge f_{j_1})\wedge\cdots\wedge(e_{i_n}\wedge f_{j_n}).
\]
A term vanishes unless the $i_k$ are pairwise distinct and the $j_k$ are
pairwise distinct, that is, unless $k\mapsto i_k$ and $k\mapsto j_k$ are
permutations $\sigma$ and $\pi$ of $\{1,\ldots,n\}$.  Since the factors
commute we may reorder them so that the primed indices appear in increasing
order; this replaces $\pi$ by $\rho=\pi\sigma^{-1}$ and leaves the coefficient
$\prod_ka_{\sigma(k)\pi(k)}=\prod_ka_{k\rho(k)}$ unchanged.  It remains to
evaluate
\[
 (e_1\wedge f_{\rho(1)})\wedge\cdots\wedge(e_n\wedge f_{\rho(n)}).
\]
Moving all the primed factors to the front costs $c_n$ by the computation of
\Cref{lem:ch09-interlaced}, and reordering
$f_{\rho(1)},\ldots,f_{\rho(n)}$ into increasing order costs
$\operatorname{sgn}(\rho)$; hence the displayed product equals
$c_n\operatorname{sgn}(\rho)\,\lambda\wedge J\lambda
=\operatorname{sgn}(\rho)\,\tau_n$, again by
\Cref{lem:ch09-interlaced}.  For each $\rho$ there are exactly $n!$ pairs
$(\sigma,\pi)$ with $\pi\sigma^{-1}=\rho$, so
\[
 \alpha^n
 =n!\sum_{\rho}\operatorname{sgn}(\rho)\prod_{k=1}^na_{k\rho(k)}\;\tau_n
 =n!\det(A)\,\tau_n .
\]

\emph{The mixed statement.}  Recall that the mixed discriminant
$D(A_1,\ldots,A_n)$ of $n$ symmetric matrices is the unique symmetric
multilinear function with
\[
 \det\Bigl(\sum_{i}t_iA_i\Bigr)
 =\sum_{i_1,\ldots,i_n}t_{i_1}\cdots t_{i_n}D(A_{i_1},\ldots,A_{i_n}),
 \qquad t_i\in\mathbf R,
\]
so that $D(A,\ldots,A)=\det A$.  Let $\alpha_1,\ldots,\alpha_n$ be
$(1,1)$-forms with coefficient matrices $A_1,\ldots,A_n$ and put
$\alpha=\sum_it_i\alpha_i$, whose coefficient matrix is $\sum_it_iA_i$.  The
map $(\alpha_1,\ldots,\alpha_n)\mapsto\alpha_1\wedge\cdots\wedge\alpha_n$ is
symmetric and multilinear, because $(1,1)$-forms commute.  Expanding
$\alpha^n$ multilinearly and comparing with
$\alpha^n=n!\det(\sum t_iA_i)\tau_n$ term by term in the monomials
$t_{i_1}\cdots t_{i_n}$ yields
\[
 \alpha_{i_1}\wedge\cdots\wedge\alpha_{i_n}
 =n!\,D(A_{i_1},\ldots,A_{i_n})\,\tau_n ,
\]
which is the assertion, with the stated normalization: on the diagonal it
returns $n!\det(A)\tau_n$.
\end{proof}

\begin{example}[Weakly positive but not positive, in the first case where
this occurs]
\label{ex:ch09-weak-not-positive}
Let $n=4$ and $p=2$, so that $S_2$ has six elements and
$\bigwedge^2V^*$ has dimension six.  For
$w=\sum_{I\in S_2}w_I\,\dprime x_I$ one computes
\[
 w\wedge w
 =2\bigl(w_{12}w_{34}-w_{13}w_{24}+w_{14}w_{23}\bigr)\lambda,
\]
because the only surviving products are those of complementary indices, and
because
\[
 \varepsilon_{12}=\varepsilon_{34}=+1,\quad
 \varepsilon_{13}=\varepsilon_{24}=-1,\quad
 \varepsilon_{14}=\varepsilon_{23}=+1 .
\]
Write $P(w)=w_{12}w_{34}-w_{13}w_{24}+w_{14}w_{23}$ for the resulting
quadratic form, the \emph{Pl\"ucker form}.  If $w=\eta_1\wedge\eta_2$ is
decomposable then $w\wedge w=0$, hence
\[
 w\in\Gamma_2\ \Longrightarrow\ P(w)=0 .
\]
Only this implication will be used.  Now let $\beta\in A^{2,2}(V)$ be the
symmetric form whose normalized coefficient matrix is
\[
 M(\beta)=\mathrm{Id}-4\widehat P,
 \qquad\text{that is,}\qquad
 w^{\mathsf T}M(\beta)w=|w|^2-4P(w),
\]
where $\widehat P$ is the symmetric matrix of $P$; explicitly,
\[
\begin{aligned}
 \beta=&-\sum_{I\in S_2}\dprime x_I\wedge\ddprime x_I
   +2\bigl(\dprime x_{12}\wedge\ddprime x_{34}
          +\dprime x_{34}\wedge\ddprime x_{12}\bigr)\\
  &-2\bigl(\dprime x_{13}\wedge\ddprime x_{24}
          +\dprime x_{24}\wedge\ddprime x_{13}\bigr)
   +2\bigl(\dprime x_{14}\wedge\ddprime x_{23}
          +\dprime x_{23}\wedge\ddprime x_{14}\bigr),
\end{aligned}
\]
the overall minus sign being $c_2=-1$.  Then:
\begin{enumerate}
\item $\beta$ is weakly positive.  Indeed, for $w\in\Gamma_2$ we have
  $P(w)=0$, hence $w^{\mathsf T}M(\beta)w=|w|^2\geq0$, and
  \Cref{prop:ch09-positivity-dictionary}\,(3) applies.
\item $\beta$ is not positive.  Take
  $w_0=\dprime x_{12}+\dprime x_{34}$, that is, the coefficient vector with
  $w_{12}=w_{34}=1$ and all other entries zero.  Then $|w_0|^2=2$ and
  $P(w_0)=1$, so
  \[
   w_0^{\mathsf T}M(\beta)w_0=2-4=-2<0,
  \]
  and $M(\beta)$ is not positive semidefinite.
\end{enumerate}
Observe that $w_0$ is not decomposable, precisely because $P(w_0)=1\neq0$;
this is the smallest instance of a nondecomposable exterior square, and it is
the entire source of the phenomenon.
\end{example}

\begin{example}[Positive but not strongly positive]
\label{ex:ch09-positive-not-strong}
Keep $n=4$, $p=2$ and $w_0=\dprime x_{12}+\dprime x_{34}$ as in
\Cref{ex:ch09-weak-not-positive}, and set
\[
 \omega=c_2\,\eta\wedge J\eta,
 \qquad
 \eta=\dprime x_1\wedge\dprime x_2+\dprime x_3\wedge\dprime x_4
     \in A^{2,0}(V),
\]
so that $M(\omega)=w_0w_0^{\mathsf T}$ by \Cref{lem:ch09-generators}.
Explicitly,
\[
 \omega=-\bigl(\dprime x_{12}+\dprime x_{34}\bigr)
        \wedge\bigl(\ddprime x_{12}+\ddprime x_{34}\bigr).
\]
By construction $\omega$ is positive; its coefficient matrix is a rank-one
square, hence positive semidefinite, in accordance with
\Cref{prop:ch09-positivity-dictionary}\,(1).  It is \emph{not} strongly
positive.  Indeed, let $\beta$ be the weakly positive form of
\Cref{ex:ch09-weak-not-positive}.  By
\Cref{lem:ch09-complement-pairing} and
$\star(w_0w_0^{\mathsf T})=(\star w_0)(\star w_0)^{\mathsf T}$, together with
$(\star w_0)_{12}=\varepsilon_{12}(w_0)_{34}=1$ and
$(\star w_0)_{34}=\varepsilon_{34}(w_0)_{12}=1$, all other entries vanishing,
we get $\star w_0=w_0$ and therefore
\[
 \omega\wedge\beta
 =\bigl\langle M(\omega),\star M(\beta)\bigr\rangle\tau_4
 =\bigl(\star w_0\bigr)^{\mathsf T}M(\beta)\bigl(\star w_0\bigr)\,\tau_4
 =w_0^{\mathsf T}M(\beta)w_0\,\tau_4=-2\,\tau_4 .
\]
If $\omega$ were strongly positive, then $\omega\wedge\beta$ would be a
nonnegative top form, by the definition of weak positivity applied to
$\beta$ and by the commutativity
$\omega\wedge\beta=\beta\wedge\omega$ of forms of even total degree.  It is
not.  Hence
\[
 \mathrm{SP}_2(\mathbf R^4)\subsetneq\mathrm{P}_2(\mathbf R^4)
 \subsetneq\mathrm{WP}_2(\mathbf R^4),
\]
the second inclusion being strict by \Cref{ex:ch09-weak-not-positive}.
\end{example}

The pair of examples repays a second reading, because the mechanism is
easily stated and easily forgotten.  All three cones are described, through
\Cref{prop:ch09-positivity-dictionary}, by one and the same matrix $M$:
strong positivity asks that $M$ be a nonnegative combination of squares
$vv^{\mathsf T}$ with $v$ \emph{decomposable}; positivity asks that $M$ be a
nonnegative combination of squares with $v$ arbitrary; weak positivity asks
only that the quadratic form $M$ be nonnegative \emph{on} decomposable
vectors.  When every vector is decomposable the three conditions collapse to
one.  When they do not, the middle condition is strictly between the other
two, and the two examples exhibit both gaps with the same six-dimensional
computation.  The same phenomenon occurs for every $n\geq4$ and every $p$
with $2\leq p\leq n-2$, by taking products with strongly positive forms in
the remaining variables; we shall not need this and do not prove it.

A last remark on terminology, since it is a standing source of confusion.
In complex geometry, as in \cite[Chapter~III]{DemBook}, only two of the
three cones are named: \emph{strongly positive} means the same there as
here, and the cone called \emph{weakly positive} here is called simply
\emph{positive} there.  The middle cone has no standard name in the complex
literature, as the dictionary below records.  It is emphatically \emph{not}
the complex \emph{strongly positive} cone: the middle cone is generated by
the forms $c_p\eta\wedge J\eta$ with $\eta\in A^{p,0}$ \emph{arbitrary},
whereas the strongly positive cone is generated by the same expressions with
$\eta$ decomposable, and \Cref{ex:ch09-positive-not-strong} separates them.

Getting this wrong makes \Cref{prop:ch09-positivity-products}\,(2) look
false.  The familiar complex statement is that the wedge of two
\emph{positive} forms need not be positive; in the names used here that is
the wedge of two \emph{weakly} positive forms, which is exactly what the
warning after that proposition says.  The middle cone, by contrast, is
closed under wedge, in the complex setting no less than here, because its
generating set is: $c_a\eta\wedge J\eta$ wedge $c_b\zeta\wedge J\zeta$ is
$c_{a+b}(\eta\wedge\zeta)\wedge J(\eta\wedge\zeta)$, again a generator,
since $\eta\wedge\zeta\in A^{a+b,0}$ and
$c_ac_b(-1)^{ab}=c_{a+b}$.  No decomposability is used, which is why the
reduction to generators is legitimate for (2) although it would not be for
a claim about the strongly positive cone.  The names used in this book are those of Lagerberg
and of \cite{CLD25}, and they are used consistently: \emph{strongly
positive} always means a nonnegative combination of the products
$\eta_1\wedge J\eta_1\wedge\cdots\wedge\eta_p\wedge J\eta_p$.  When a
statement is imported from the complex literature, the translation is made
explicitly.

\begin{example}[Rank and vanishing of top powers]
\label{ex:ch09-rank-one}
For a linear functional $\ell$ and a one-variable function $g$,
\Cref{lem:superform-composition} gives
\[
 \dprimeddprime(g\circ\ell)=g''(\ell)\,
                 \dprime\ell\wedge\ddprime\ell.
\]
Its matrix has rank at most one, so its square vanishes.  When $g''\geq0$
the form is positive.  This elementary rank calculation is the local model
for Monge--Amp\`ere measures supported on lower-dimensional tropical strata.
\end{example}

\section{Convexity as local plurisubharmonicity}
\label{sec:ch09-convexity}

We come to the statement for whose sake the formalism was built.  In complex
pluripotential theory a psh function is one whose $\ddc$ is a positive
$(1,1)$-current; in the Lagerberg calculus the corresponding class is
described by an entirely classical condition, and the whole of Part~II may be
read as the systematic exploitation of the following identification.

Recall the notation: for $f\in C^2(U)$ we write $f_i=\partial f/\partial x_i$
and $f_{ij}=\partial^2f/\partial x_i\partial x_j$, and
$\operatorname{Hess}f=(f_{ij})$.

\begin{theorem}
\label{thm:convex-ddc-positive}
Assume that $U$ is convex. For $f\in C^2(U)$ the following are equivalent:
\begin{enumerate}
\item $f$ is convex;
\item the Hessian matrix $(f_{ij})$ is positive semidefinite;
\item $\dprimeddprime f$ is a positive $(1,1)$-form.
\end{enumerate}
\par
\end{theorem}
\begin{proof}
The equivalence of (1) and (2) follows by restricting $f$ to affine lines:
the second derivative along $a+tv$ is $v^T(\operatorname{Hess}f)v$.
The displayed formula for $\dprimeddprime f$ identifies its coefficient matrix with
the Hessian. The definition of positivity in bidegree $(1,1)$ is exactly
positive semidefiniteness of that matrix.

In detail.  Suppose (1).  Fix $a\in U$ and $v\in V$ and let
$I=\{t\in\mathbf R:a+tv\in U\}$, an open interval since $U$ is convex and
open.  The function $g(t)=f(a+tv)$ is convex on $I$, being the restriction of
a convex function to a segment, and it is $C^2$; hence $g''\geq0$.  By the
chain rule,
\[
 g''(t)=\sum_{i,j}f_{ij}(a+tv)v_iv_j
 =v^{\mathsf T}\operatorname{Hess}f(a+tv)\,v .
\]
Taking $t=0$ and letting $a$ and $v$ vary gives (2).  Conversely, assume (2)
and let $a,b\in U$.  Applying the one-variable Taylor formula with integral
remainder to $g(t)=f\bigl(a+t(b-a)\bigr)$ on $[0,1]$, which is legitimate
because the segment lies in $U$,
\[
 f(b)=f(a)+\sum_if_i(a)(b_i-a_i)
   +\int_0^1(1-t)\,g''(t)\,\rd t
 \geq f(a)+\langle\nabla f(a),b-a\rangle .
\]
Thus $f$ lies above each of its tangent affine functions, which for a
differentiable function is equivalent to convexity: given
$b_0,b_1\in U$ and $s\in[0,1]$, apply the inequality at
$a=(1-s)b_0+sb_1$ to $b=b_0$ and to $b=b_1$ and combine the two with weights
$1-s$ and $s$; the gradient terms cancel and one obtains
$f(a)\leq(1-s)f(b_0)+sf(b_1)$.  This proves (1).

The equivalence of (2) and (3) is now a matter of reading
\Cref{ex:ch09-hessian-form}: the normalized coefficient matrix of
$\dprimeddprime f$ is $\operatorname{Hess}f$, since $c_1=1$, and by
\Cref{prop:superform-matrix-wedge} a $(1,1)$-form is positive exactly when
that matrix is positive semidefinite at every point.  By
\Cref{prop:ch09-positivity-inclusions} the three positivity notions agree in
bidegree $(1,1)$, so it is immaterial which of them is meant in (3).
\end{proof}

If $U$ is not convex, the same proof gives a local statement: positivity of
$\dprimeddprime f$ is equivalent to convexity on every sufficiently small
convex open subset of $U$.  One should not define global convexity on a
nonconvex domain by comparing values across a line segment that leaves the
domain.

This is the first appearance of a distinction which will recur throughout
Part~II.  Positivity of $\dprimeddprime f$ is a \emph{local} condition: it is
checked at each point separately, and it is preserved under restriction to
open subsets.  Convexity is a \emph{global} condition on the domain: it
compares values at pairs of points, and it is meaningful only when the domain
contains the segment joining them.  On a convex domain the two agree.  On a
general domain the correct statement is that positivity of
$\dprimeddprime f$ is equivalent to \emph{local convexity}, that is, to
convexity in a neighbourhood of each point.  The passage from local to global
is exactly the content of the maximum principles of
\Cref{chap:cld-bedford-taylor}, and it is not automatic; on a tropical space
with nontrivial topology it can fail.

\begin{warning}{Convexity is not preserved by nonaffine coordinate changes}
Positivity of $\dprimeddprime f$ is preserved by affine pullback, by
\Cref{prop:ch09-positivity-inclusions} and
\Cref{prop:ch09-pullback-functorial}, and by nothing more.  If
$F:\mathbf R\to\mathbf R$ is a nonaffine diffeomorphism, then $f\circ F$
need not be convex when $f$ is: take $f(x)=x$, which is convex, and
$F(x)=-x^2$ on $x>0$, so that $f\circ F=-x^2$ is concave.  This is not a
defect of the calculus; it is the reason that the transition maps of the
tropical charts in \Cref{chap:tropical-spaces} are required to be integral
affine, and the reason that a psh function on a Berkovich space cannot be
defined by pulling back convex functions along arbitrary maps.
\end{warning}

\begin{corollary}
\label{cor:superform-convex-composition}
Let $f$ be smooth and convex on a convex open set, and let $\chi$ be smooth,
convex, and nondecreasing on an interval containing $f(U)$. Then
$\chi\circ f$ is convex and $\dprimeddprime(\chi\circ f)$ is positive.
\end{corollary}
\begin{proof}
In \Cref{lem:superform-composition}, both coefficients are nonnegative,
$\dprimeddprime f$ is positive by \Cref{thm:convex-ddc-positive}, and
$\dprime f\wedge\ddprime f$ is strongly positive.  The conclusion follows
from stability of the positive cone under addition and multiplication by
nonnegative functions.
\end{proof}

\begin{example}[Log-sum-exp]
For affine functions $\ell_1,\ldots,\ell_N$ and $\varepsilon>0$, put
\[
 L_\varepsilon(x)=\varepsilon\log
       \left(\sum_{k=1}^N\rme^{\ell_k(x)/\varepsilon}\right).
\]
If
\[
 p_k(x)=\frac{\rme^{\ell_k(x)/\varepsilon}}
               {\sum_j\rme^{\ell_j(x)/\varepsilon}},
\]
then the Hessian is $\varepsilon^{-1}$ times the covariance matrix of the
linear parts of the $\ell_k$ with probabilities $p_k(x)$.  It is therefore
positive semidefinite.  Moreover,
\[
 \max_k\ell_k\leq L_\varepsilon
 \leq\max_k\ell_k+\varepsilon\log N.
\]
Thus maxima of finitely many affine functions admit explicit smooth convex
approximants.  The same regularization is used later to pass from real PL psh
functions to smooth CLD-psh functions.
\end{example}

\subsection{The Monge--Amp\`ere form in coordinates}

\begin{proposition}
\label{prop:ch09-ma-coordinates}
Let $u,u_1,\ldots,u_n\in C^2(U)$.  Then
\[
 \bigl(\dprimeddprime u\bigr)^n
 =n!\,\det\bigl(\operatorname{Hess}u\bigr)\,\tau_n,
\]
\[
 \dprimeddprime u_1\wedge\cdots\wedge\dprimeddprime u_n
 =n!\,D\bigl(\operatorname{Hess}u_1,\ldots,
             \operatorname{Hess}u_n\bigr)\,\tau_n,
\]
where $D$ is the mixed discriminant.  If $u$ is convex, then
$(\dprimeddprime u)^k$ is a strongly positive $(k,k)$-form for
$0\leq k\leq n$; in particular $\det(\operatorname{Hess}u)\geq0$.  If
$u_1,\ldots,u_n$ are convex, then
$D(\operatorname{Hess}u_1,\ldots,\operatorname{Hess}u_n)\geq0$.
\end{proposition}

\begin{proof}
The two displayed identities are
\Cref{prop:superform-matrix-wedge} applied to the $(1,1)$-forms
$\dprimeddprime u_i$, whose coefficient matrices are the Hessians by
\Cref{ex:ch09-hessian-form}.  If $u$ is convex then
$\dprimeddprime u$ is positive by \Cref{thm:convex-ddc-positive}, hence
strongly positive because the three cones coincide in bidegree $(1,1)$
(\Cref{prop:ch09-positivity-inclusions}); so
$(\dprimeddprime u)^k$ is strongly positive by
\Cref{prop:ch09-positivity-products}\,(1).  For $k=n$ a strongly positive
top form is a nonnegative multiple of $\tau_n$, whence
$n!\det(\operatorname{Hess}u)\geq0$.  The mixed statement follows in the same
way, the product of $n$ strongly positive $(1,1)$-forms being a strongly
positive top form.
\end{proof}

The factor $n!$ is not removable and is not a matter of taste.  It is the
same factor that appears in the mass ledger of
\Cref{sec:ch08-dict-three-column} and in the identity
$\tfrac1{n!}(\dprimeddprime u)^n=\det(D^2u)\,\rd x$ recorded in
\Cref{sec:appD-supercurrents}: with the integration convention fixed in
\Cref{chap:tropical-integration}, for which $\int f\,\tau_n=\int f\,\rd x$
against the covolume-one Lebesgue measure, the top wedge power of
$\dprimeddprime u$ has total mass $n!$ times the Alexandrov measure of
\Cref{thm:appD-alexandrov-measure}.  Whenever a Monge--Amp\`ere measure is
compared with a geometric intersection number, that factor must be visible;
the convention of this book is never to hide it.

\begin{example}[Monge--Amp\`ere of a convex function in two variables]
\label{ex:ch09-ma-plane}
Let $n=2$.

\emph{A quadratic potential.}  For
$u(x,y)=\tfrac12\bigl(ax^2+2bxy+cy^2\bigr)$ one has
\[
 \dprimeddprime u
 =a\,\dprime x\wedge\ddprime x
  +b\bigl(\dprime x\wedge\ddprime y+\dprime y\wedge\ddprime x\bigr)
  +c\,\dprime y\wedge\ddprime y,
\]
and by \Cref{ex:ch09-plane},
\[
 \bigl(\dprimeddprime u\bigr)^2=2\bigl(ac-b^2\bigr)\,\tau_2 .
\]
The function is convex exactly when $a\geq0$, $c\geq0$ and $ac\geq b^2$; the
Monge--Amp\`ere form is then a nonnegative constant multiple of $\tau_2$, as
\Cref{prop:ch09-ma-coordinates} requires.  Note that $ac-b^2\geq0$ alone does
\emph{not} imply convexity: the matrix $-\mathrm{Id}$ has determinant $1$.
Positivity of a $(1,1)$-form is positive semidefiniteness, not
nonnegativity of the top power.

\emph{A degenerate potential.}  Let
\[
 u(x,y)=\log\bigl(\rme^{x}+\rme^{y}\bigr),
\]
which is the case $\varepsilon=1$, $N=2$, $\ell_1=x$, $\ell_2=y$ of the
log-sum-exp example above, hence convex.  Writing
$p=\rme^{x}/(\rme^{x}+\rme^{y})$ and $q=1-p$, a direct computation gives
\[
 u_{xx}=pq,\qquad u_{xy}=u_{yx}=-pq,\qquad u_{yy}=pq,
\]
so that
\[
 \dprimeddprime u
 =pq\,\bigl(\dprime x-\dprime y\bigr)\wedge
        \bigl(\ddprime x-\ddprime y\bigr)
 =pq\;\dprime\ell\wedge\ddprime\ell,
 \qquad \ell=x-y,
\]
a positive form of rank one, and
\[
 \bigl(\dprimeddprime u\bigr)^2=2\bigl((pq)^2-(pq)^2\bigr)\tau_2=0 .
\]
The Monge--Amp\`ere form vanishes identically.  This is the analytic shadow of
an elementary fact: $u(x,y)=x+\log(1+\rme^{y-x})$ is an affine function plus a
function of the single variable $y-x$, so its Hessian has rank at most one
everywhere, exactly as in \Cref{ex:ch09-rank-one}.

\emph{Where the mass goes.}  As $\varepsilon\downarrow0$ the functions
$\varepsilon\log(\rme^{x/\varepsilon}+\rme^{y/\varepsilon})$ decrease to
$\max\{x,y\}$, whose Monge--Amp\`ere measure in the sense of
\Cref{thm:appD-alexandrov-measure} is again zero, the subgradient image
being the segment joining $(1,0)$ and $(0,1)$, of planar measure zero.  The
picture changes for $\max\{0,x,y\}$: there the subgradient image at the
origin is the full triangle with vertices $(0,0)$, $(1,0)$, $(0,1)$, of area
$\tfrac12$, and the Monge--Amp\`ere measure is an atom, as computed in
\Cref{ex:appD-polyhedral-atom}.  No smooth superform can be that measure,
which is why \Cref{chap:cld-currents} is needed; the corresponding
smooth-form statement is simply that
$(\dprimeddprime u_\varepsilon)^2$ concentrates near the origin as
$\varepsilon\downarrow0$ while keeping total mass $2!\cdot\tfrac12=1$.
\end{example}

\section{Local regularization and compactness tests}
\label{sec:ch09-regularization}

Fix an auxiliary Euclidean norm on $V$ and a nonnegative smooth function
$\rho$ supported in the unit ball with integral one.  Write
$\rho_\varepsilon(x)=\varepsilon^{-n}\rho(x/\varepsilon)$.  If
$K\Subset U$ and $\varepsilon<\operatorname{dist}(K,V\setminus U)$, the
convolution
\[
                       f_\varepsilon=f*\rho_\varepsilon
\]
is defined near $K$.

\begin{proposition}
\label{prop:superform-convex-mollification}
If $f$ is continuous and convex on a convex open set $U$, then
$f_\varepsilon$ is smooth and convex wherever it is defined, and
$f_\varepsilon\to f$ locally uniformly.  No monotonicity in $\varepsilon$ is
asserted.
\end{proposition}
\begin{proof}
For fixed $y$, the translate $x\mapsto f(x-y)$ is convex.  A positive average
of convex functions is convex, so $f_\varepsilon$ is convex.  Smoothness is
the standard differentiation-under-convolution argument.  On a compact
subset of $U$, continuity of a convex function is uniform; hence
\[
 |f_\varepsilon(x)-f(x)|
 \leq\sup_{|y|\leq\varepsilon}|f(x-y)-f(x)|
\]
tends uniformly to zero.  The final warning is essential: ordinary
mollifications need not form a decreasing family.
\end{proof}

This elementary proposition isolates two recurring themes.  Positivity is
stable under averaging, but a monotone approximation requires an additional
construction.  Also, the Euclidean norm used for convolution is auxiliary;
the limiting notion of convexity and the superform calculus are affine
invariant.

The additional construction is not elaborate: it is enough to choose the
kernel radial.  The resulting statement is used so often in
\Cref{chap:cld-bedford-taylor} and in \Cref{part:bj} --- every time a
Bedford--Taylor product is defined by regularization, and every time a
monotone family is needed to justify an interchange of limits --- that it is
worth proving here once, with all four assertions.

\begin{proposition}
\label{prop:ch09-radial-mollification}
Assume in addition that $\rho$ is \emph{radial}, that is,
$\rho(y)$ depends only on $|y|$; in particular $\rho(-y)=\rho(y)$.  Let $U$
be convex and open, let $f:U\to\mathbf R$ be convex, and for $\varepsilon>0$
put
\[
 U_\varepsilon=\{x\in U:\ x+\varepsilon\overline B\subset U\},
\]
a convex open set, with $U_\varepsilon\uparrow U$ as
$\varepsilon\downarrow0$.  Then, on $U_\varepsilon$:
\begin{enumerate}
\item $f_\varepsilon=f*\rho_\varepsilon$ is smooth and convex;
\item $f_\varepsilon\geq f$;
\item if $0<\varepsilon\leq\varepsilon'$ then
  $f_\varepsilon\leq f_{\varepsilon'}$ on $U_{\varepsilon'}$;
\item $f_\varepsilon\downarrow f$ pointwise as $\varepsilon\downarrow0$, and
  the convergence is locally uniform.
\end{enumerate}
\end{proposition}

\begin{proof}
That $U_\varepsilon$ is convex follows from the description
$U_\varepsilon=\bigcap_{|y|\leq\varepsilon}(U-y)$ as an intersection of
convex sets, and it is open because $U$ is open and $\overline B$ compact.
Assertion (1) is \Cref{prop:superform-convex-mollification}.

Fix $x\in U_{\varepsilon_0}$ and define, for $|t|\leq\varepsilon_0$,
\[
 g_x(t)=\int_Vf(x+ty)\,\rho(y)\,\rd y .
\]
The integrand is defined because $x+ty\in U$ whenever $|t|\leq\varepsilon_0$
and $y\in\operatorname{supp}\rho\subset\overline B$.  Substituting
$y\mapsto-y$ and using $\rho(-y)=\rho(y)$ gives
\[
 f_\varepsilon(x)=\int_Vf(x-\varepsilon y)\rho(y)\,\rd y=g_x(\varepsilon),
 \qquad 0\leq\varepsilon\leq\varepsilon_0,
\]
and in particular $g_x(0)=f(x)$.  Two properties of $g_x$ are immediate.  It
is convex, being an average with the nonnegative weights $\rho(y)\rd y$ of
the convex functions $t\mapsto f(x+ty)$.  And it is even:
\[
 g_x(-t)=\int f(x-ty)\rho(y)\rd y=\int f(x+ty)\rho(-y)\rd y=g_x(t).
\]

A convex even function on $[-\varepsilon_0,\varepsilon_0]$ is nondecreasing
on $[0,\varepsilon_0]$.  Indeed, for $0\leq s<t\leq\varepsilon_0$ write
\[
 s=\frac{t-s}{2t}\,(-t)+\frac{t+s}{2t}\,(t),
\]
a convex combination since the two weights are nonnegative and sum to $1$;
convexity and evenness give
\[
 g_x(s)\leq\frac{t-s}{2t}g_x(-t)+\frac{t+s}{2t}g_x(t)=g_x(t).
\]
Taking $s=0$ and $t=\varepsilon$ yields
$f(x)=g_x(0)\leq g_x(\varepsilon)=f_\varepsilon(x)$, which is (2); taking
$s=\varepsilon$ and $t=\varepsilon'$ yields (3), valid at every
$x\in U_{\varepsilon'}$.

Finally, (4).  By (3) the family $(f_\varepsilon)$ decreases pointwise as
$\varepsilon\downarrow0$, and by
\Cref{prop:superform-convex-mollification} it converges to $f$ locally
uniformly; a decreasing family converging in any sense to $f$ converges to
$f$ pointwise.
\end{proof}

\begin{example}[Both extra conclusions fail for a general kernel]
\label{ex:ch09-nonradial-kernel}
Radiality is not a convenience.  Take $n=1$, let $\rho\geq0$ be smooth with
integral one and support contained in the interval $(0,1)$, and put
$m=\int y\rho(y)\rd y\in(0,1)$.  For the affine, hence convex, function
$f(x)=x$ one computes
\[
 f_\varepsilon(x)=\int(x-\varepsilon y)\rho(y)\,\rd y=x-\varepsilon m .
\]
Thus $f_\varepsilon<f$ everywhere, so (2) of
\Cref{prop:ch09-radial-mollification} fails, and
$\varepsilon\mapsto f_\varepsilon$ is strictly \emph{decreasing}, so (3)
fails as well.  Inspecting the proof shows exactly which hypothesis is
responsible for which conclusion: the domination $f_\varepsilon\geq f$ needs
only that the kernel have barycentre $0$, that is
$\int y\rho(y)\rd y=0$, which is Jensen's inequality; the monotonicity in
$\varepsilon$ needs the evenness $\rho(-y)=\rho(y)$.  A radial kernel has
both properties, and this is the only reason to insist on one.
\end{example}

\subsection{The distributional statement}

The functions to which the theory is applied---piecewise affine potentials,
maxima of finitely many affine functions, limits of model
metrics---are almost never of class $C^2$.  What is needed is therefore the
distributional form of \Cref{thm:convex-ddc-positive}.  We state it on an
open subset of $\mathbf R^n$, in the language of distributions rather than of
currents: the current-theoretic version, which is the one used from
\Cref{chap:cld-currents} onwards, has the same proof once the duality
formalism is available.

For $u\in L^1_{\loc}(U)$ the coefficients of $\dprimeddprime u$ are the
distributional second derivatives $\partial_i\partial_ju$, and \emph{$\dprimeddprime u\geq0$
in the sense of distributions} means, in accordance with
\Cref{thm:convex-ddc-positive}\,(3) and
\Cref{prop:ch09-positivity-dictionary}, that
\begin{equation}
\label{eq:ch09-distributional-positivity}
 \sum_{i,j}v_iv_j\int_Uu\,\frac{\partial^2\chi}{\partial x_i\partial x_j}
   \,\rd x\;\geq\;0
 \qquad\text{for all }v\in\mathbf R^n
 \text{ and all }\chi\in C_c^\infty(U),\ \chi\geq0 .
\end{equation}

\begin{theorem}
\label{thm:ch09-distributional-convexity}
Let $U\subset\mathbf R^n$ be open and convex and let $u\in L^1_{\loc}(U)$.
The following are equivalent:
\begin{enumerate}
\item $\dprimeddprime u\geq0$ in the sense of
  \eqref{eq:ch09-distributional-positivity};
\item there is a convex function $\tilde u:U\to\mathbf R$ with
  $\tilde u=u$ almost everywhere.
\end{enumerate}
The function $\tilde u$ is then unique, and, for a radial kernel $\rho$ as in
\Cref{prop:ch09-radial-mollification},
\[
 \tilde u(x)=\lim_{\varepsilon\downarrow0}\,u*\rho_\varepsilon(x)
 \qquad\text{for every }x\in U,
\]
the family $u*\rho_\varepsilon$ being convex and nonincreasing as
$\varepsilon\downarrow0$.
\end{theorem}

\begin{proof}
Throughout, $\rho$ is radial and $u_\varepsilon=u*\rho_\varepsilon$, defined
and smooth on $U_\varepsilon$.  We use two standard properties of mollifiers,
for which we refer to \cite{EvansGariepy}: $u_\varepsilon\to u$ in
$L^1_{\loc}(U)$ as $\varepsilon\downarrow0$, and $u_\varepsilon\to u$
locally uniformly if $u$ is continuous.

\emph{(2) $\Rightarrow$ (1).}  A finite convex function on an open convex
subset of $\mathbf R^n$ is locally Lipschitz, hence continuous and locally
integrable; see \cite[\S10]{Roc70}.  Replacing $u$ by $\tilde u$ changes
neither side of \eqref{eq:ch09-distributional-positivity}, so we may assume
$u=\tilde u$ is convex.  By
\Cref{prop:ch09-radial-mollification}\,(1) each $u_\varepsilon$ is smooth and
convex on $U_\varepsilon$, so its Hessian is positive semidefinite by
\Cref{thm:convex-ddc-positive}.  Let $\chi\in C^\infty_c(U)$ be nonnegative
and $v\in\mathbf R^n$.  For $\varepsilon$ small the support of $\chi$ lies in
$U_\varepsilon$, and integrating by parts twice,
\[
 \sum_{i,j}v_iv_j\int u_\varepsilon\,\partial_i\partial_j\chi\,\rd x
 =\int\Bigl(\sum_{i,j}v_iv_j\partial_i\partial_ju_\varepsilon\Bigr)
   \chi\,\rd x\geq0 .
\]
Letting $\varepsilon\downarrow0$ and using $u_\varepsilon\to u$ in
$L^1_{\loc}$ gives \eqref{eq:ch09-distributional-positivity}.

\emph{(1) $\Rightarrow$ (2).}  Fix $\varepsilon>0$ and $x\in U_\varepsilon$.
The function $y\mapsto\rho_\varepsilon(x-y)$ is nonnegative, smooth and
compactly supported in $U$, so \eqref{eq:ch09-distributional-positivity}
applies to it.  Since
$\partial^2_{v,y}\bigl[\rho_\varepsilon(x-y)\bigr]
=(\partial_v^2\rho_\varepsilon)(x-y)$, differentiating under the integral
sign gives
\[
 \sum_{i,j}v_iv_j\,\partial_i\partial_ju_\varepsilon(x)
 =\int_Uu(y)\,(\partial_v^2\rho_\varepsilon)(x-y)\,\rd y\;\geq\;0 .
\]
Hence $u_\varepsilon$ has positive semidefinite Hessian on the convex open
set $U_\varepsilon$, and is convex there by
\Cref{thm:convex-ddc-positive}.

Next we show that $\varepsilon\mapsto u_\varepsilon(x)$ is nondecreasing.
Let $0<\varepsilon\leq\varepsilon'$ and let $x$ be a point at which both are
defined, say $x\in U_{\varepsilon'+\delta_0}$ for some $\delta_0>0$.  For
$0<\delta\leq\delta_0$ the function $u_\delta$ is convex on $U_\delta$, so
\Cref{prop:ch09-radial-mollification}\,(3) applied to it gives
\[
 (u_\delta)_{\varepsilon}(x)\leq(u_\delta)_{\varepsilon'}(x).
\]
Convolution is commutative and associative, so
$(u_\delta)_{t}=u*\rho_\delta*\rho_{t}=(u_{t})_\delta$ for $t\in\{\varepsilon,\varepsilon'\}$.
As $\delta\downarrow0$ we have $(u_t)_\delta(x)\to u_t(x)$, because $u_t$ is
continuous.  Passing to the limit in the displayed inequality yields
$u_\varepsilon(x)\leq u_{\varepsilon'}(x)$.

Therefore the limit
$\tilde u(x)=\lim_{\varepsilon\downarrow0}u_\varepsilon(x)$ exists in
$[-\infty,+\infty)$ for every $x\in U$, the family being nonincreasing as
$\varepsilon\downarrow0$ and defined for all small $\varepsilon$.  Being a
pointwise decreasing limit of convex functions, $\tilde u$ satisfies the
convexity inequality
$\tilde u((1-s)x_0+sx_1)\leq(1-s)\tilde u(x_0)+s\tilde u(x_1)$, all terms
being limits of decreasing real sequences and hence well defined in
$[-\infty,+\infty)$.

We claim $\tilde u$ is finite.  Since $u_\varepsilon\to u$ in
$L^1_{\loc}$, some sequence $\varepsilon_k\downarrow0$ satisfies
$u_{\varepsilon_k}\to u$ almost everywhere; combined with the pointwise
monotone convergence this gives $\tilde u=u$ almost everywhere, so
$\tilde u>-\infty$ outside a null set.  Suppose $\tilde u(x_0)=-\infty$ for
some $x_0\in U$ and choose $r>0$ with $B(x_0,r)\subset U$.  The reflection
$y\mapsto2x_0-y$ preserves $B(x_0,r)$ and Lebesgue measure, so the set
$Z=\{y\in B(x_0,r):\tilde u(y)=-\infty\}$, which is null, has null
reflection; hence there is $y\in B(x_0,r)$ with $\tilde u(y)$ and
$\tilde u(2x_0-y)$ both finite.  The convexity inequality at the midpoint
$x_0$ then gives
$\tilde u(x_0)\leq\tfrac12(\tilde u(y)+\tilde u(2x_0-y))>-\infty$, a
contradiction.  Thus $\tilde u$ is a finite convex function equal to $u$
almost everywhere, which is (2), and the displayed formula for $\tilde u$ has
been proved along the way.

\emph{Uniqueness.}  Two finite convex functions on an open convex set which
agree almost everywhere are equal, both being continuous by
\cite[\S10]{Roc70}.
\end{proof}

\begin{corollary}
\label{cor:ch09-monotone-smooth-approximation}
Let $U$ be convex and open and let $u:U\to\mathbf R$ be convex.  Then there
exists a family $(u_\varepsilon)_{\varepsilon>0}$ of smooth convex functions,
each defined on $U_\varepsilon$, such that
\[
 u_\varepsilon\downarrow u\quad\text{pointwise and locally uniformly as }
 \varepsilon\downarrow0,
 \qquad
 \dprimeddprime u_\varepsilon\geq0 .
\]
In particular every convex function on a convex open subset of $V$ is a
decreasing limit of smooth functions with positive $\dprimeddprime$.
\end{corollary}

\begin{proof}
Take $u_\varepsilon=u*\rho_\varepsilon$ with $\rho$ radial and apply
\Cref{prop:ch09-radial-mollification} together with
\Cref{thm:convex-ddc-positive}.
\end{proof}

\Cref{cor:ch09-monotone-smooth-approximation} is the local prototype of every
approximation theorem in this book.  It should be compared with the two
places where the analogous statement is hard or false.  On a Berkovich space
the psh functions of \Cref{chap:cld-bedford-taylor} are \emph{defined} by an
approximation property, precisely because no convolution is available;
and in the global theories of \Cref{part:bfj,part:bj} the approximating
objects are model functions, whose existence is a theorem about models rather
than an integral formula.  What survives from the present section is the
shape of the statement: a decreasing family of regular objects, converging
locally uniformly, each of which is positive.

\begin{remark}[The boundary of the smooth category]
The function $\max\{0,x\}$ is convex but not smooth; its second derivative is
a Dirac mass. Hence curvature of a piecewise affine model metric is generally
not a smooth superform. \Cref{chap:cld-currents} treats it as a current, while
\Cref{chap:delta-forms} introduces delta-forms, where such polyhedral
curvature can be multiplied.
Harmonic tropicalizations enlarge the chart category for a different reason:
they repair cohomological limitations of the original smooth sheaf. These two
extensions solve different problems and will not be conflated.
\end{remark}

\section{Dictionary and sign checklist}
\label{sec:ch09-dictionary}

\subsection{The three-column dictionary}

The following table is the translation key between the three languages that
meet in this chapter.  It is not a theorem; each row abbreviates a statement
proved above or, in the middle column, a statement of complex geometry that
is recalled for orientation only.

\begin{center}
\small
\begin{tabular}{p{.30\textwidth}|p{.30\textwidth}|p{.30\textwidth}}
\textbf{Superform calculus}&\textbf{Complex analogue}
  &\textbf{Convex-analytic content}\\ \hline
$A^{p,q}(U)$, bidegree $(p,q)$
  &forms of type $(p,q)$
  &---\\
$\dprime,\ \ddprime$
  &$\partial,\ \bar\partial$
  &one and the same real derivative, twice labelled\\
$J$, $J^2=\operatorname{id}$
  &complex conjugation
  &transposition of the coefficient array\\
$\dprimeddprime u$
  &$\ddc u$, up to \Cref{app:normalizations}
  &the Hessian $\operatorname{Hess}u$\\
symmetric $(p,p)$-form
  &real $(p,p)$-form
  &symmetric matrix indexed by $S_p$\\
strongly positive
  &strongly positive
  &$\sum_sa_sv_sv_s^{\mathsf T}$, $v_s$ decomposable\\
positive
  &self-dual cone, no standard name
  &$M$ positive semidefinite\\
weakly positive
  &positive, in \cite[Ch.~III]{DemBook}
  &$w^{\mathsf T}Mw\geq0$ for $w$ decomposable\\
$\tau_n$ positive
  &positive volume form
  &nonnegative density\\
$\dprimeddprime u\geq0$
  &$u$ plurisubharmonic
  &$u$ convex\\
$\tfrac1{n!}(\dprimeddprime u)^n$
  &Monge--Amp\`ere measure
  &$\det(D^2u)\,\rd x$, the Alexandrov measure\\
$F^*$ for $F$ affine
  &pullback by a holomorphic map
  &affine change of variables\\
$\iota'_v,\ \iota''_v$
  &contraction by a $(1,0)$- or $(0,1)$-vector
  &directional derivative after
   \Cref{prop:ch09-cartan}
\end{tabular}
\end{center}

\subsection{The sign checklist}

Seven signs govern everything above.  Each has been proved; each is worth
memorizing in the form of the test that detects an error.

\begin{enumerate}
\item \emph{The constants.}  $c_p=(-1)^{p(p-1)/2}$ and
  $c_ac_b(-1)^{ab}=c_{a+b}$
  (\Cref{lem:ch09-interlaced}).  \emph{Test:} $c_0=c_1=1$, $c_2=c_3=-1$,
  $c_4=1$.
\item \emph{The orientation.}
  $\tau_n=c_n\,\dprime x_1\wedge\cdots\wedge\dprime x_n\wedge
   \ddprime x_1\wedge\cdots\wedge\ddprime x_n$
  (\Cref{lem:ch09-interlaced}).  \emph{Test:} for $n=2$ the grouped product
  is $-\tau_2$.
\item \emph{The block sign.}  Moving a double-primed block of degree $q$
  past a primed block of degree $p'$ costs $(-1)^{qp'}$
  (\Cref{prop:ch09-wedge-signs}).  \emph{Test:} the commutation of two
  forms depends only on total degrees.
\item \emph{The operator sign.}  $\ddprime$ carries $(-1)^p$ on $A^{p,q}$,
  and so does the graded contraction $\widetilde\iota''_v=(-1)^p\iota''_v$;
  $\dprime$, $\iota'_v$ and the block contraction $\iota''_v$ carry none
  (\Cref{def:ch09-operators}, \Cref{def:ch09-contraction}).  \emph{Test:}
  $\ddprime=J\dprime J$ and $\widetilde\iota''_v=J\iota'_vJ$.  It is the
  block contraction that appears in the boundary traces
  \eqref{eq:superform-boundary-traces}, which is why those carry an explicit
  $(-1)^{n-1}$.
\item \emph{The involution.}  Here $J\ddprime x=\dprime x$ and symmetry is
  $J\alpha=(-1)^p\alpha$; in \cite{CLD25} the involution differs by
  $(-1)^q$ on $A^{p,q}$ and symmetry reads $J\alpha=\alpha$
  (\Cref{rem:ch09-two-involutions}).  \emph{Test:} both conventions agree on
  $A^{p,0}$, so the positivity cones are the same.
\item \emph{The Hessian normalization.}
  $\dprimeddprime u=\sum_{i,j}u_{ij}\dprime x_i\wedge\ddprime x_j$, with no
  factor $\tfrac12$, $2\pi$, or $\rmi$
  (\Cref{ex:ch09-hessian-form} and \eqref{eq:appJ-real-hessian}).
  \emph{Test:} on $\mathbf R$, $\dprimeddprime u=u''\tau_1$.
\item \emph{The determinant.}  $(\dprimeddprime u)^n=n!\det(\operatorname{Hess}u)\tau_n$
  and $F^*\tau_n=(\det B)^2\tau_n$
  (\Cref{prop:ch09-ma-coordinates,prop:ch09-pullback-top}).  \emph{Test:}
  the factor $n!$ is present in every mass computation, and the square in
  every change of variables.
\end{enumerate}

\subsection*{Statements exported to later chapters}

The following are the results of this chapter that are quoted elsewhere, with
their first use.
\begin{itemize}
\item The sign rules
  \Cref{prop:ch09-wedge-signs,lem:ch09-interlaced} and the differential
  identities \Cref{prop:superform-differential-identities}: used throughout
  \Cref{chap:tropical-spaces,chap:tropical-integration,chap:cld-forms}.
\item The contraction calculus
  \Cref{def:ch09-contraction,lem:ch09-contraction-basic,prop:ch09-cartan}:
  used for the boundary traces \eqref{eq:superform-boundary-traces} in
  \Cref{chap:tropical-integration} and for the homogeneity conditions
  \eqref{eq:appD-strong-homogeneity}.
\item Affine functoriality
  \Cref{prop:ch09-pullback-functorial,prop:ch09-pullback-top}: the reason
  tropical transition maps are integral affine
  (\Cref{chap:tropical-spaces}) and the source of the lattice factors in
  \Cref{chap:tropical-integration}.
\item The positivity dictionary and duality
  \Cref{prop:ch09-positivity-dictionary,thm:ch09-positivity-duality} and the
  product rules \Cref{prop:ch09-positivity-products}: used for positive
  currents in \Cref{chap:cld-currents} and for the Bedford--Taylor products
  of \Cref{chap:cld-bedford-taylor}.
\item The convexity criteria
  \Cref{thm:convex-ddc-positive,thm:ch09-distributional-convexity} and the
  regularization results
  \Cref{prop:ch09-radial-mollification,cor:ch09-monotone-smooth-approximation}:
  used in \Cref{chap:cld-bedford-taylor} and, in their global counterparts,
  in \Cref{part:bfj,part:bj}.
\item The Monge--Amp\`ere normalization
  \Cref{prop:ch09-ma-coordinates}: the factor $n!$ that reappears in every
  comparison between a Monge--Amp\`ere mass and an intersection number.
\end{itemize}

\section{Source notes}

Superforms were introduced by Lagerberg. Our signs and positivity cones agree
with the 2025 revision of Chambert--Loir--Ducros
\cite[\S\S1--2]{CLD25}. The presentation begins with the Hessian test because
this is the portion used most frequently in pluripotential theory.

The bigraded algebra $A^{\bullet,\bullet}$, the operators $\dprime$ and
$\ddprime$, the involution, the three positivity cones, and the associated
supercurrents are due to A.~Lagerberg, in his paper on super currents and
tropical geometry (\emph{Math.\ Z.}, 2012).  That work is not in the
bibliography of this book, because no proof here takes an unproved statement
from it; it is named rather than cited, and the reader who wishes to consult
it will find its results reproved above with the conventions of this book.
Two of Lagerberg's contributions should be attributed precisely: the
observation that the doubled differential calculus reproduces the formal
apparatus of complex geometry over $\mathbf R$, and the definition of the
three positivity cones, which are the real analogues of the cones of
\cite[Ch.~III]{DemBook}.  The example of a nonzero weakly positive
$(2,2)$-form on $\mathbf R^4$ orthogonal to every strongly positive form is
also his; our
\Cref{ex:ch09-weak-not-positive,ex:ch09-positive-not-strong} are a variant,
arranged so that both strictness assertions follow from a single
six-dimensional computation.

The development of that algebra into a theory of forms and currents on
Berkovich spaces is due to Chambert-Loir and Ducros
\cite{CLD12}, in the revised form \cite{CLD25}.  The distinction matters at
three points in this chapter.  First, the involution: \cite{CLD25} normalizes
$J$ so that $J^2=(-1)^{p+q}$, an almost-complex-structure convention, whereas
we keep Lagerberg's involution with $J^2=\operatorname{id}$; the translation
is \Cref{rem:ch09-two-involutions} and it affects no positivity statement.
Second, the terminology: what is called \emph{weakly positive} by Lagerberg
and by \cite{CLD25}, and here, is what most of the complex literature calls
positive; the warning after
\Cref{ex:ch09-positive-not-strong} records this.  Third, the ordering:
\cite{CLD25} develops the pointwise algebra, then differential superforms,
then integration; we have inserted the contraction operators and the
coefficient-matrix dictionary early, because they make the positivity theory
a matter of symmetric matrices, and we have postponed integration entirely to
\Cref{chap:tropical-integration}, where the noninvariance of the naive
Lagerberg integral is the motivation for calibrations.

The linear algebra of the positivity cones is classical in its complex form;
\cite[Ch.~III]{DemBook} is the standard reference for the complex cones,
their duality, and the failure of the inclusions to be equalities.  The
convex-analytic inputs---the bipolar theorem for closed convex cones, the
continuity of finite convex functions on open sets, and the subgradient
description of the Monge--Amp\`ere measure---are quoted from \cite{Roc70},
and the convex-body material that reappears in
\Cref{chap:tropical-spaces} may be found in \cite{Sch14}.  Standard
properties of mollifiers are used as in \cite{EvansGariepy}.
The finite-dimensional statements we import are collected, with their exact
forms, in \Cref{app:convexity}.

Two later reworkings of the same formalism are worth naming here, since they
change the emphasis rather than the content.  Berndtsson
\cite{Ber19Super,Ber26Convex} develops superforms and supercurrents with a
Riemannian and convex-geometric emphasis, obtaining in particular the
supercurrent formulation of mixed volumes used in
\Cref{sec:appD-supercurrents}; his metric-dependent supercurrent of a real
submanifold is a genuinely different object from the calibrated tropical
currents of \Cref{chap:tropical-integration}, and the two are not identified
in this book.  Gubler and K\"unnemann \cite{GK17,GK19} enlarge the category
of forms so that polyhedral curvature can be multiplied, and Gubler, Jell and
Rabinoff \cite{GJR26} enlarge the category of charts to repair the
cohomological behaviour of the smooth sheaf; a Poincar\'e lemma for the
associated real forms is \cite{Jell16}.  These extensions solve different
problems and are treated separately in \Cref{chap:delta-forms}.  The
tropical-geometric background for the polyhedral chapters is
\cite{MR18Tropical}.

Finally, a word on what this chapter deliberately does not do.  It does not
define superforms with coefficients in a subsheaf, a refinement that
\cite{CLD25} uses systematically and that is needed for piecewise linear and
$G$-rational structures; that material is developed where it is used, in
\Cref{chap:tropical-spaces} and \Cref{app:tropical-algebra}.  It does not
treat integration, currents, or any singular object.  And it does not assert
that a pointwise conic decomposition of a smooth positive form can be chosen
smoothly; the pointwise convention adopted in
\Cref{sec:ch09-positivity} is the one used throughout the book, and no
argument anywhere depends on the stronger reading.
